% Part E fragments owned by cycles agent. No preamble.
\subsection{E.1: the missing saturation condition on 14:9}
\label{sec:effectivity-saturation}

\paragraph{E.1a (Proved: saturation as a precise polarity).}
Let $G/\mathbb Q$ be a reductive algebraic group, with its category of
finite-dimensional rational representations. Let $\mathcal I$ be the
set of complex cocharacters, modulo conjugation if desired. For
$C\subset\mathcal I$ and a class $\mathcal E$ of representations,
define
\[
\begin{split}
 E(C)&=\{V:\text{every weight of every }\nu\in C
                 \text{ on }V_{\mathbb C}\text{ is nonnegative}\},\\
 I(\mathcal E)&=\{\nu:\text{every weight of }\nu
                 \text{ on every }V\in\mathcal E
                 \text{ is nonnegative}\},\\
 \operatorname{Sat}(C)&=I(E(C)).
\end{split}
\]
Then $\operatorname{Sat}$ is extensive, monotone, and idempotent, and
$E(\operatorname{Sat}(C))=E(C)$. The exact missing condition is
\begin{equation}\label{eq:effective-saturation}
 C=I(E(C)).
\end{equation}
The saturated cocharacter sets are in inclusion-reversing bijection
with the full representation subcategories that are intersections of
the $P(\nu)$. Additional positivity, Tate-twist, weight-zero, and
Galois-invariance axioms are imposed on the objects on both sides;
they do not replace \eqref{eq:effective-saturation}.

\emph{Proof.}
The defining condition gives the Galois connection
$\mathcal E\subset E(C)$ if and only if $C\subset I(\mathcal E)$.
In particular $C\subset I(E(C))$. Applying the order-reversing map
$E$ gives $E(I(E(C)))\subset E(C)$, while the reverse inclusion is
true from the definition of $I$. Thus $E(I(E(C)))=E(C)$, which also
proves idempotence. If $\mathcal E=\bigcap_{\nu\in C}P(\nu)=E(C)$,
its full set of effective cocharacters is $I(\mathcal E)$ and is
saturated. Conversely a saturated $C$ is recovered from the
intersection defining $E(C)$. This proves the bijection. The same
argument applies to pro-reductive groups by using their actual
finite-dimensional representations; it does not require choosing
one finite-dimensional maximal torus of a pro-group.

\paragraph{E.1b (Proved lattice form).}
For a torus with character lattice $M$ and cocharacter lattice
$N=\operatorname{Hom}(M,\mathbb Z)$, assume $C\subset N$ is Galois
stable. Then
\[
 C^\vee=\{m\in M:\langle m,c\rangle\ge0\quad\forall c\in C\},
 \qquad
 \operatorname{Sat}(C)=
 \{v\in N:\langle m,v\rangle\ge0\quad\forall m\in C^\vee\}.
\]
To check that rational representations detect every listed character,
we take the representation whose weights are the finite Galois orbit
of $m$. Galois stability of $C$ makes the whole orbit effective
exactly when $m$ is in $C^\vee$. For a split connected reductive group,
the same formula works on a maximal torus for a Weyl-stable cone:
move $m$ into a dominant chamber; the irreducible representation of
highest weight $m$ has all its weights in the convex hull of its Weyl
orbit. The Weyl-stable polar therefore makes that whole representation
effective. Over $\mathbb Q$ we use a rational representation containing
the finite Galois orbit, with multiplicity if descent requires it;
multiplicities do not change inequalities.

\paragraph{E.1c (Refuted automaticity for additive cones).}
We give complete data compatible with the normalization
$\langle\varepsilon,j\rangle=2$. Let
\[
 T=\mathbb G_m\times R^1_{\mathbb Q(i)/\mathbb Q}\mathbb G_m,
 \qquad N=\mathbb Z^2,
 \qquad \overline{(a,b)}=(a,-b),
 \quad \varepsilon=(1,0)\in M,
 \quad j=(2,0)\in N.
\]
Set
\[
 C=\mathbb N(1,1)+\mathbb N(1,-1)
   =\{(a,b)\in\mathbb Z^2:a\ge |b|, a\equiv b\pmod2\}.
\]
It is an additive Galois-stable cone, contains $j$, and has
$\varepsilon$ strictly positive away from the origin.
The character $\varepsilon$ is in the interior of its real polar.
Its integral polar is
$C^\vee=\{(u,v):u\ge|v|\}$, because the two generator inequalities
are $u+v\ge0$ and $u-v\ge0$. Therefore
\[
 \operatorname{Sat}(C)=\{(a,b):a\ge|b|\}\supsetneqq C;
 \qquad (1,0)\in\operatorname{Sat}(C)\setminus C.
\]
The two cones define exactly the same effectivity structure. This
is a counterexample to automatic saturation under the additive-cone
requirements alone.

Even imposing root saturation is not enough. The monoid
\[
 C_{\mathrm{open}}=\{0\}\cup
     \{(a,b)\in\mathbb Z^2:a>|b|\}
\]
is root saturated: $nv\in C_{\mathrm{open}}$ for an integer $n>0$
implies $v\in C_{\mathrm{open}}$. It has the same integral polar and
the same bipolar as $C$, but omits $(1,1)$ and $(1,-1)$. Thus the
required saturation includes closure as detected by integral
half-spaces, not merely divisibility in the cocharacter lattice.

\paragraph{E.1d (Proved automaticity under a stronger meaning of ``cone'').}
There is an important qualification to the preceding answer. Suppose
``cone'' is defined in advance to mean
$C=K\cap N$, where $K\subset N_{\mathbb R}$ is a closed convex cone
and $\varepsilon$ lies in the interior of $K^\vee$.
Then $C$ is saturated as above, in particular when $K$ is rational
polyhedral. The conclusion holds without rational polyhedrality:
a separating
normal for a lattice point outside $K$ can be perturbed into the
interior of $K^\vee$ while keeping its value negative at that point,
then approximated by a rational normal and scaled to an integral
character. This character is nonnegative on $K\cap N$ and excludes
the point from the integral bipolar.

For completeness, the perturbation uses
$m_t=m+t\varepsilon$ for sufficiently small $t>0$, where
$m\in K^\vee$ satisfies $\langle m,v\rangle<0$.
Since $\varepsilon$ is interior, $m_t$ is interior; a sufficiently
close rational point remains there and remains negative on $v$.
The separating-hyperplane theorem supplies $m$ because $K$ is closed
and convex. This proves both inclusions of the integral bipolar.

Thus the manuscript's unspecified quotation-marked ``c\^one'' matters.
Saturation is not a consequence of the stated additive requirements;
it is automatic if the stronger closed-convex-lattice definition is
already part of those requirements. Formula
\eqref{eq:effective-saturation} removes that ambiguity entirely.
It also gives the strongest exact classification requested in E.1.

\paragraph{E.1e (Audit of the portion of R13 used here).}
The closure of $P(\nu)$ under sums, subobjects and tensor products
follows by taking weights, and a positive
$\langle\varepsilon,\nu\rangle$ makes sufficiently many Tate twists
of every fixed representation effective. For a cone with the strict
interior condition, the same twist works uniformly on that cone:
normalize $\langle\varepsilon,c\rangle=1$, use compactness of the
closed section, and bound each of the finitely many weights of the
representation on it. This verifies the step of R13 needed in E.1.
We do not use E.1 as a proof of the motivic conjecture that every
cohomologically effective motive is effective. Lead E1, concerning
the Galois orbit and Frobenius slopes, belongs to E.2 and is not
decided by these purely representation-theoretic statements.

\subsection{E.3: Albanese equivalence and the Abel--Jacobi obstruction}
\label{sec:albanese-family-audit}

\paragraph{E.3a (Proved definitions and the part of R14 used).}
Let $N^i(X)\subset\mathrm{CH}^i(X)_{\mathrm{alg}}$ be generated by
$W_*\nu$, where $T$ is smooth projective connected,
$W\in\mathrm{CH}^i(X\times T)$ is a family and
$\nu\in\mathrm{CH}_0(T)^0$ has Albanese class zero. Set
\[
 P^i(X)=\mathrm{CH}^i(X)_{\mathrm{alg}}/N^i(X).
\]
The graph of a correspondence carries such generators to generators,
so this construction is functorial for correspondences. For divisors,
the map from a parameter variety to $\operatorname{Pic}^0(X)$ factors
through its Albanese, hence $N^1(X)=0$ in $\mathrm{CH}^1(X)$.
For zero-cycles, the diagonal proves that the whole Albanese kernel
is in $N^n(X)$; conversely functoriality of the Albanese map shows
that every generator has zero Albanese class. Thus
\[
 P^1(X)=\operatorname{Pic}^0(X)(k),\qquad
 P^n(X)=\operatorname{Alb}(X)(k).
\]
These are the two directions of the manuscript's Theorem 1.4, rather
than merely its stated conclusion.

The product of two algebraically trivial classes is in $N^*(X)$:
write them as fibre differences over parameter varieties $T,T'$;
their product is obtained from the rectangle on $T\times T'$.
Both coordinate sums of that rectangle vanish, so its Albanese
class vanishes. This proves the portion of Theorem 1.6 used here.

The replacement of a parameter variety by an abelian parameter,
used below, is valid up to an integer. Choose a smooth curve
$C\subset T$ whose Jacobian surjects onto $\operatorname{Alb}(T)$.
Poincar\'e complete reducibility supplies a homomorphism
$s:\operatorname{Alb}(T)\to J(C)$ with
$q s=[m]$ for some $m>0$. The normalized Poincar\'e divisor on
$J(C)\times C$, followed by $C\to T$, gives a family $U$ with
$\operatorname{alb}(U(a))=ma$. Composing $U$ with $W$ realizes
$m$ times the original map by an abelian family. The equality
$U(a)=m\nu$ is an equality modulo Albanese equivalence when
$a=\operatorname{alb}(\nu)$, not an equality of cycles. This verifies
the missing construction used by R14 on 14:32--34.

\paragraph{E.3b (Proved in characteristic zero: Theorem 3.1).}
Let $X/\mathbb C$ have dimension $n$, let $A/\mathbb C$ be an
abelian variety, and let $Z\in\mathrm{CH}^i(X\times A)$, normalized
by its fibre at zero. It defines a group homomorphism
$u_Z:A(\mathbb C)\to P^i(X)$ and a cohomological action
\[
 a({}^tZ):H^{2(n-i)+1}(X,\mathbb Q)(n-i)
              \longrightarrow H^1(A,\mathbb Q).
\]
If $u_Z=0$, then $a({}^tZ)=0$.

\emph{Proof.}
The Abel--Jacobi map on algebraically trivial cycles is regular,
with values in the algebraic intermediate Jacobian
$J_a^i(X)$; see the opened
\citet[Example 1.14, Theorem 9.3]{cyACMV2022}. For any generator
$W_*\nu$ of $N^i(X)$, the regular map
$t\mapsto\operatorname{AJ}(W_t-W_{t_0})$ to this abelian variety
factors through $\operatorname{Alb}(T)$. Hence its value on $\nu$
is zero. We obtain a well-defined homomorphism
\[
 \overline{\operatorname{AJ}}:P^i(X)\longrightarrow J_a^i(X)(\mathbb C).
\]
The normalized family $Z$ gives a morphism of abelian varieties
$\phi_Z:A\to J_a^i(X)$ satisfying
$\phi_Z(a)=\overline{\operatorname{AJ}}(u_Z(a))$.
Thus $u_Z=0$ implies $\phi_Z=0$.

Its map on $H_1$ is the cylinder map of the family. To see this
without a degree convention hidden in notation, transport $Z_a$
along a closed path in $A$. The swept-out chain is a cycle of real
dimension $2(n-i)+1$ in $X$. Integrating a closed form over that
cycle gives the period change of the Abel--Jacobi integral. This is
exactly the pushforward of the K\"unneth component of $[Z]$ with
one leg in $H^1(A)$; dualizing gives the displayed
$a({}^tZ)$. Therefore $\phi_Z=0$ makes that cohomological map zero.
Comparison gives the same result with $\mathbb Q_\ell$ coefficients.

For an arbitrary algebraically closed field of characteristic zero,
one uses a finitely generated field of definition for the cycle
and the descended regular Abel--Jacobi homomorphism
\citep[Theorems 9.1 and 9.3]{cyACMV2022}. If its cohomological map
were nonzero, the induced homomorphism of the descended abelian
varieties would be nonzero, hence nonzero on a point over an
algebraic closure of that finitely generated field. This point and
any proposed Albanese-equivalence relation spread over finitely
generated characteristic-zero fields and embed in $\mathbb C$,
contradicting the complex case. This explains the direction of
transfer; simply embedding a possibly very large algebraically
closed ground field in $\mathbb C$ would not be valid.

\paragraph{E.3c (Proved exact rational reduction of the converse: Lead E4).}
For fixed $X/\mathbb C$, the following are equivalent:
\begin{enumerate}
\item For every abelian family $Z$, the equality
$a({}^tZ)=0$ implies $u_Z\otimes\mathbb Q=0$.
\item The homomorphism
$\overline{\operatorname{AJ}}\otimes\mathbb Q$
is injective on the portion of $P^i(X)\otimes\mathbb Q$ generated
by abelian families.
\end{enumerate}
By E.3a this portion is all of $P^i(X)\otimes\mathbb Q$.
Thus the rational formulation of Lead E4 is correct. Its unqualified
integral formulation is not established by this argument;
E.3h--k below give a counterexample to that stronger biconditional.

\emph{Proof.}
Assume (2). If $a({}^tZ)=0$, the homomorphism of abelian varieties
$\phi_Z$ has zero rational $H_1$ action and is zero. Hence
$u_Z(a)$ is in the Abel--Jacobi kernel for every $a$, and (2)
makes its rational class zero.

Conversely let $z$ have zero rational Abel--Jacobi image. Using
E.3a and multiplying by a positive integer, we write
$z=u_Z(a)$ for one abelian family (a finite sum is combined using
a product of parameter varieties). After another multiplication
we have $\phi_Z(a)=0$. The kernel of $\phi_Z$ has finitely many
components. Multiplying $a$ by their exponent puts it in
$K=(\ker\phi_Z)^0$, an abelian subvariety of $A$.
Restrict $Z$ to $X\times K$. Its homomorphism to $J_a^i(X)$ is
zero, hence its cohomological action is zero by the same
cylinder-map description. Hypothesis (1) makes its $u$ zero
after tensoring with $\mathbb Q$, so the original $z$ is zero
there as well. This proves (2).

The integer multiplications in this proof are essential to its
stated scope. It does not prove integral equality of the two
equivalence relations, nor does the phrase ``cohomologically zero''
silently discard finite components of a kernel.

\paragraph{E.3d (Proved endpoint cases; exact remaining gates).}
For $i=1$ and $i=n$, both directions follow in every characteristic
from $P^1=\operatorname{Pic}^0(k)$ and
$P^n=\operatorname{Alb}(k)$: the family map is an actual homomorphism
of abelian varieties, and its rational Tate-module action is zero
if and only if the homomorphism is zero. The inverse direction
uses faithfulness, not the Tate conjecture over a finitely generated
field.

For intermediate codimension, E.3c is a reduction, not a proof of
its injectivity statement. A first attempted argument in characteristic
$p$ tried to apply the torsion Bloch map separately to each generator
$W_*\nu$ of $N^i(X)$. That argument is invalid because those
generators need not be torsion. E.3e below resolves this obstruction
by specializing the entire torsion relation at once. The opened
\citet[Theorem 8.1, Proposition 8.2]{cyACMV2022} provides stronger
results in codimension two for algebraic representatives, but
does not by itself identify the manuscript's quotient $P^i(X)$
with that representative.

\paragraph{E.3e (Proved: Theorem 3.1 in every positive characteristic).}
Let $k$ be algebraically closed of characteristic $p>0$, let
$\ell\ne p$, let $X/k$ be smooth projective of dimension $n$, and
use projective smooth parameter varieties in the definition of
Albanese equivalence. For an abelian variety $A/k$ and a normalized
cycle $Z\in\mathrm{CH}^i(A\times X)$, we have
\[
 u_Z=0\quad\Longrightarrow\quad
 a({}^tZ):H^{2(n-i)+1}(X,\mathbb Q_\ell(n-i))
                       \longrightarrow H^1(A,\mathbb Q_\ell)
 \text{ is zero}.
\]
For integral $Z$ the corresponding map with $\mathbb Z_\ell$
coefficients is also zero, since its target $H^1(A,\mathbb Z_\ell)$
is torsion free. Thus the characteristic-$p$ implication requested
on 14:38 is true.
The proof does not require an algebraic adjoint in the middle
cohomology of $X$.

\emph{The torsion lemma.}
Bloch's map
\[
 \lambda_X^i:\mathrm{CH}^i(X)[\ell^\infty]
             \longrightarrow H^{2i-1}(X,\mathbb Q_\ell/\mathbb Z_\ell(i))
\]
vanishes on
$N^i(X)\cap\mathrm{CH}^i(X)[\ell^\infty]$.
We prove this carefully, without requiring the individual
Albanese-equivalence generators to be torsion.

Take $z$ in that intersection. We choose an expression
$z=\sum_j (W_j)_*\nu_j$, where $T_j$ are smooth projective,
$\nu_j$ are degree-zero zero-cycles on $T_j$, and
$\operatorname{alb}_{T_j}(\nu_j)=0$. We also choose rational
equivalences expressing this equality and $\ell^r z=0$ for some
$r$. All cycles, functions, morphisms and points in these finite
data descend to a field $K$ finitely generated over $\mathbb F_p$.
We include the Albanese varieties $B_j$ of the geometric $T_j$
and their pointed Albanese maps $b_j:T_j\to B_j$ in the descended
data. After enlarging $K$ and shrinking an integral finite-type
model over $\mathbb F_p$, these extend to smooth projective
families, abelian schemes, cycles, and morphisms. We shrink further
so that the supports and rational functions used in the finite
Chow relations spread and the fibrewise relations remain valid.
In particular, we preserve the chosen relation $\ell^r z=0$.
The equality
$b_{j*}\nu_j=0$ holds as an equality of sections of $B_j$.

Choose a closed point of this model in the good locus. The
specialization can be taken along a valuation with geometric
generic field $\overline K$ and residue field
$\overline{\mathbb F}_p$: at a smooth closed point, a full flag
of local parameters gives an iterated discrete valuation, and
we extend it to $\overline K$. Equivalently we specialize along
that flag one step at a time. Proper specialization of Chow
classes commutes with the proper pushforwards and refined
pullbacks defining $(W_j)_*\nu_j$. We therefore obtain
\[
 z_s=\sum_j(W_{j,s})_*\nu_{j,s},
 \qquad \ell^r z_s=0,\qquad b_{j,s*}\nu_{j,s}=0.
\]

We must justify that specialization does not silently change
the needed Albanese kernel. The map
$b_j^*:H^1(B_j,\mathbb Z_\ell)\to H^1(T_j,\mathbb Z_\ell)$ is an
isomorphism on the geometric generic fibre, integrally: it is
the usual identification with the dual of the Albanese Tate
module. This is a morphism of lisse $\mathbb Z_\ell$ sheaves
on our connected good model. Smooth proper base change
therefore makes it an isomorphism at $s$ as well. Thus the induced
homomorphism
$\operatorname{Alb}(T_{j,s})\to B_{j,s}$ is an isogeny with
degree prime to $\ell$. The Albanese class of $\nu_{j,s}$ lies
in its finite kernel and hence has order prime to $\ell$.
Every degree-zero Chow class over $\overline{\mathbb F}_p$ is
torsion, as proved in C.5a by curve parameterization. Roitman's
prime-to-$p$ torsion theorem, in the proved form
\citet[Theorem 4.2]{cyBloch1979}, identifies its $\ell$-primary
torsion with that of the Albanese. Consequently $\nu_{j,s}$ has
no $\ell$-primary part: it has finite order prime to $\ell$.
So $z_s$ has order prime to $\ell$ while also being killed by
$\ell^r$, and therefore $z_s=0$.

The Bloch map commutes with specialization
\citep[Proposition 3.8]{cyBloch1979}, and smooth proper base
change identifies the corresponding cohomology groups with
$\mathbb Q_\ell/\mathbb Z_\ell$ coefficients. Hence
$\lambda_X^i(z)=0$. Extension from $\overline K$ back to $k$
preserves this conclusion by the same base-change compatibility.
This proves the torsion lemma.

\emph{Application to an abelian family.}
Put $g=\dim A$ and view $Z$ as a correspondence from $A$ to $X$.
For every $\alpha\in\mathrm{CH}^g(A)[\ell^\infty]$, the
definition of $u_Z$ gives
$Z_*\alpha\in N^i(X)$ if $u_Z=0$; moreover $Z_*\alpha$ remains
$\ell$-primary torsion. The lemma and correspondence
functoriality of the Bloch map
\citep[Proposition 3.5]{cyBloch1979} imply
$\lambda_X^i Z_*\alpha=0$. On zero-cycles the map
$\lambda_A^g$ identifies with the torsion Albanese map and is
an isomorphism
\citep[Proposition 3.9 and Theorem 4.2]{cyBloch1979}.
Taking Tate modules first gives zero on the torsion-free
$\mathbb Z_\ell$ cohomology quotients. Here we use
$T_\ell H^m(-,\mathbb Q_\ell/\mathbb Z_\ell(r))
=H^m(-,\mathbb Z_\ell(r))/\mathrm{tors}$; the possible finite
torsion in the next cohomology group contributes no Tate module.
Tensoring with $\mathbb Q_\ell$ therefore shows that
\[
 Z_*:H^{2g-1}(A,\mathbb Q_\ell(g))
           \longrightarrow H^{2i-1}(X,\mathbb Q_\ell(i))
\]
is zero. Its Poincar\'e dual is exactly the displayed
$a({}^tZ)$, so the theorem follows. Finally, for integral $Z$,
the integral transpose map takes values in the free module
$H^1(A,\mathbb Z_\ell)$. Its rationalization is zero, so the
integral map itself is zero. We do not assert torsion
injectivity for the Bloch map in codimension $i$.

\paragraph{E.3f (Refuted: replacing the integral premise by its rationalization).}
In characteristic $p$, the premise in E.3e cannot be weakened to
$u_Z\otimes\mathbb Q=0$. Take
$k=\overline{\mathbb F}_5$, let $A=X$ be the elliptic curve
$y^2=x^3-x$ with origin at infinity, and put
$Z=\Delta_X-X\times\{0\}$. Its discriminant is
$64\equiv4\pmod5$, so it is smooth. Here
$P^1(X)=\operatorname{Pic}^0(X)(k)\simeq X(k)$ and $u_Z$ is
the identity on this nonzero torsion group. Thus
$u_Z\otimes\mathbb Q=0$, whereas the action on $H^1(X,\mathbb Q_\ell)$
is the identity for every $\ell\ne5$ and is nonzero. This example
explains why the proof retains integral torsion cycles until
the final cohomological step.

% Further E.3 deductions; owned by cycles. No preamble.
\paragraph{E.3g (Proved: integral abelian parameterization and regular invariants).}
Every algebraically trivial cycle $z\in\mathrm{CH}^i(X)_{\mathrm{alg}}$
on a smooth projective complex variety is, \emph{integrally}, the
fibre difference at a point and at zero of a cycle on $X\times A$
for some abelian variety $A$. Moreover every regular homomorphism
\[
 \rho:\mathrm{CH}^i(X)_{\mathrm{alg}}\longrightarrow B(\mathbb C)
\]
to an abelian variety kills $N^i(X)$.

For the first assertion, the definition of algebraic equivalence
expresses $z$ as a sum of $\Gamma_{j*}\delta_j$, with $C_j$ smooth
projective connected curves, $\delta_j\in\mathrm{CH}_0(C_j)^0$,
and $\Gamma_j\in\mathrm{CH}^i(C_j\times X)$. A normalized
Poincar\'e divisor on $J(C_j)\times C_j$ has fibre the divisor
class corresponding to any point of $J(C_j)=\operatorname{Pic}^0(C_j)$.
Composing this divisor correspondence with $\Gamma_j$ therefore
realizes $\Gamma_{j*}\delta_j$ exactly. Pulling the finitely many
families to $A=\prod_jJ(C_j)$ and adding them realizes $z$.
No isogeny denominator is involved. The up-to-an-integer argument
in E.3a remains necessary for its stronger assertion about
replacing a whole prescribed parameter variety by its own Albanese.

For the second assertion, regularity makes
$t\mapsto\rho(W_t-W_{t_0})$ a morphism from each smooth projective
connected parameter variety $T$ to $B$, taking $t_0$ to zero.
It therefore factors through $\operatorname{Alb}(T)$. Its value on
each $\nu\in\ker(\mathrm{CH}_0(T)^0\to\operatorname{Alb}(T))$
is zero. These images generate $N^i(X)$, proving the assertion.

The group $N^i(X)$ is divisible over $\mathbb C$. Indeed
$\mathrm{CH}_0(T)^0$ is divisible by curve parameterization. If
$\nu$ is in its Albanese kernel and $mw=\nu$, then
$\operatorname{alb}(w)\in\operatorname{Alb}(T)[m]$.
Roitman's torsion isomorphism
\citep[Theorem 4.2]{cyBloch1979} supplies
$\tau\in\mathrm{CH}_0(T)^0[m]$ with the same Albanese class.
Thus $w-\tau$ is in the Albanese kernel and $m(w-\tau)=\nu$.
Images and sums of divisible groups remain divisible, proving
the claim for $N^i(X)$. This is compatible with the independent
fact that the Walker kernel is divisible, whereas the ordinary
Abel--Jacobi kernel need not be
\citep[Theorem 2.3 and Remark 2.4, author p.~12]{cyACMVWalker2022}.

In particular the surjective regular Walker map
\[
 \psi_W:\mathrm{CH}^i(X)_{\mathrm{alg}}\longrightarrow J_W^i(X),
 \qquad
 J_W^i(X)=J\bigl(N^{i-1}H^{2i-1}(X,\mathbb Z(i))\bigr),
\]
factors through $P^i(X)$. Here $N^{i-1}$ in the cohomology group
is the coniveau filtration, not the subgroup $N^i(X)$ defining
Albanese equivalence. The regularity and surjectivity are the
opened \citet[Theorem 1.1, pp.~2--3; Corollaries 2.5--2.6,
pp.~7--8]{cySuzuki2021}. Thus we have
\[
 N^i(X)\subset\ker\psi_W\subset\ker\operatorname{AJ},
 \qquad
 \operatorname{AJ}=\alpha\psi_W,
\]
where $\alpha:J_W^i(X)\to J_a^i(X)$ is the natural isogeny.

\paragraph{E.3h (Refuted: integral equality with ordinary Abel--Jacobi equivalence).}
There exists a smooth projective complex fourfold $X$ and an
algebraically trivial codimension-three cycle $z$ such that
\[
 \operatorname{AJ}(z)=0,
 \qquad z\notin N^3(X).
\]
By E.3g this $z$ comes from an abelian family integrally. Hence the
integral equality of equivalence relations asserted on the right
side of Lead E4 is false. The same failure occurs for some
$d$-fold in every codimension $3\le i\le d-1$, $d\ge4$.

We use the exact opened input of
\citet[Corollary 4.2 and Remark 4.3, author p.~10]{cyOttemSuzuki2020}.
They construct a smooth threefold $W$, a rank-one degeneracy locus
of a map
\[
 \mathcal O^{\oplus3}\longrightarrow
 \mathcal O(1,2,0)\oplus\mathcal O(1,0,2)
 \quad\hbox{on }\mathbb P^1\times\mathbb P^2\times\mathbb P^2,
\]
and an elliptic curve $E$, for which $X=W\times E$ has a
surjective regular Walker map and the isogeny
$\alpha:J_W^3(X)\to J_a^3(X)$ has nonzero finite kernel.
More precisely, their construction takes coordinates $[S:T]$
on the first factor and the rank-one locus of the matrix
\[
 \begin{pmatrix}
 P_1&Q_1&R_1\\
 SP_2+qP_3&SQ_2+qQ_3&SR_2+qR_3
 \end{pmatrix}.
\]
Here $q$ is a sufficiently large auxiliary prime, and the three
triples of rational-coefficient polynomials have tridegrees
$(1,2,0)$, $(0,0,2)$ and $(1,0,2)$, respectively, with the
generality and specialization conditions imposed in their
Theorem 3.1. These are the exact construction data in that
theorem, not a selected numerical example.
Their Theorem 3.1 supplies $W$ over $\mathbb Q$ with torsion-free
integral cohomology and a nonalgebraic integral Hodge class; their
Theorem 4.1 supplies $E$. We use this as a proved primary theorem,
not as a claim that we have independently verified their full
degeneration argument or selected numerical coefficients for $W,E$.
E.3r below gives a fully displayed geometric generic matrix and
Legendre elliptic curve to which this obstruction transfers
integrally, and supplies a nonzero cycle of order two.

Choose $0\ne b\in\ker\alpha$ and use surjectivity to choose
$z$ with $\psi_W(z)=b$. Then $\operatorname{AJ}(z)=0$, while
E.3g gives $N^3(X)\subset\ker\psi_W$, so $z\notin N^3(X)$.
This is the complete deduction from the stated primary input.
For $d\ge4$ and $3\le i\le d-1$, put $s=d-4$, $j=i-3$ and
take $z'=z\times t^j$ on $X\times\mathbb P^s$. Its Abel--Jacobi
image is zero by functoriality. If $z'$ belonged to $N^i$, then
multiplication by $t^{s-j}$ followed by projection would put $z$
in $N^3(X)$, by functoriality of Albanese equivalence. Since
$0\le j\le s$, the projection is exactly $z$, a contradiction.

\paragraph{E.3i (Proved qualification: the finite obstruction does not settle
the family-map converse).}
The counterexample in E.3h does not itself give an abelian family
with $a({}^tZ)=0$ but $u_Z\ne0$. For any abelian family, regularity
gives an actual homomorphism
$\phi_{W,Z}:A\to J_W^i(X)$ with
$\alpha\phi_{W,Z}=\phi_Z$. If $a({}^tZ)=0$, then
$\phi_Z=0$ by the cylinder-map argument in E.3b. Thus the image
of $\phi_{W,Z}$ is contained in the finite group $\ker\alpha$.
An abelian variety is connected, so this homomorphism is zero.
Consequently the Walker invariant cannot detect a failure of the
family-map converse. The individual point $z$ in E.3h is realized
by a family whose whole Abel--Jacobi homomorphism need not vanish.

This identifies a precise logical defect in the unqualified
integral version of Lead E4: passage from a point in the kernel
of an abelian-variety homomorphism to its connected kernel requires
an integer multiple. E.3c proves the rational equivalence, and
E.3h refutes the proposed integral equality of equivalence relations.
We do not infer that the integral family-map converse is false.
E.3j--k below supply the additional argument needed to decide the
integral biconditional on this same fourfold.

\paragraph{E.3j (Proved: the family converse for one-cycles when zero-cycles
are supported on a curve).}
Let $X/\mathbb C$ be a smooth projective connected $n$-fold,
$n\ge2$, and suppose that a closed subvariety $C\subset X$
of dimension at most one supports $\mathrm{CH}_0(X)$; that is,
$\mathrm{CH}_0(C)\to\mathrm{CH}_0(X)$ is surjective. Then for
every abelian variety $A$ and every normalized
$Z\in\mathrm{CH}^{n-1}(A\times X)$,
\[
 a({}^tZ)=0\quad\Longrightarrow\quad u_Z=0
 \quad\text{in }\operatorname{Hom}\bigl(A(\mathbb C),P^{n-1}(X)\bigr).
\]
This is an integral conclusion about $u_Z$, although the
cohomological vanishing can be tested with rational coefficients.

\emph{The decomposition and its precise cycle action.}
We give the decomposition-of-the-diagonal argument needed here.
Choose a countable algebraically closed subfield $k_0\subset\mathbb C$
over which $X$ and $C$ are defined. Embed $K=k_0(X)$ in $\mathbb C$
over $k_0$. The corresponding geometric generic point $\eta$,
viewed as a complex point of $X$, is rationally equivalent to a
zero-cycle supported on $C$. The finitely many cycles and rational
functions in this relation spread over a finite-type integral
$\overline K$-scheme. A closed $\overline K$-point in a good open
specializes the relation to $\overline K$, preserving $\eta$ and
the support on $C$. The resulting finite data descend to a finite
extension $K'/K$. Taking the norm therefore gives
\[
 M[\eta]=z_C\quad\text{in }\mathrm{CH}_0(X_K),
 \qquad M>0,\quad\operatorname{supp}(z_C)\subset C_K.
\]
Closing $z_C$ in $X\times C$ and applying localization to the
generic point of the first factor yields
\begin{equation}\label{eq:E3-support-decomposition}
 M\Delta_X=\Gamma_C+\Gamma_D\quad\text{in }\mathrm{CH}^n(X\times X),
\end{equation}
where $\Gamma_C$ is supported on $X\times C$ and $\Gamma_D$ on
$D\times X$ for a divisor $D\subset X$. We have proved the
needed decomposition from the stated complex support hypothesis;
we have not assumed universal support over every field extension.

Normalize the curve components of $C$. Their contribution lifts
to divisor correspondences on $X\times\widetilde C$; components
over isolated points have the same, or simpler, factorization.
They act on $\mathrm{CH}^{n-1}(X)$ through
$\mathrm{CH}^{0}(\widetilde C)$, since the lifted correspondence
has codimension one and the source has dimension $n$.
As $\mathrm{CH}^{0}(\widetilde C)$ has no algebraically trivial
part, $\Gamma_C$ annihilates
$\mathrm{CH}^{n-1}(X)_{\mathrm{alg}}$.

For each component of $D$ choose a projective resolution
$f_j:T_j\to D_j\hookrightarrow X$, with $T_j$ smooth of
dimension $n-1$. After rational coefficients, every cycle
supported on $D\times X$ lifts to a sum of cycles on
$T_j\times X$. To check surjectivity even over singular strata,
take a component of the inverse image of the given cycle that
dominates it, then a projective multisection of its generic
fibre. The closure is generically finite of positive degree;
division by that degree gives the desired lift. Clearing the
finitely many denominators in \eqref{eq:E3-support-decomposition},
we obtain an integer, still denoted $M>0$, and integral cycles
$G_j\in\mathrm{CH}^{n-1}(T_j\times X)$ such that, on
algebraically trivial cycles,
\[
 M\operatorname{id}
       =\sum_j(G_j)_*f_j^*.
\]
All these operations preserve Albanese equivalence by E.3a.
Thus the same identity factors through
\begin{equation}\label{eq:E3-Alb-factorization}
 P^{n-1}(X)\xrightarrow{(f_j^*)_j}
       \prod_j\operatorname{Alb}(T_j)(\mathbb C)
       \xrightarrow{\sum_j(G_j)_*}P^{n-1}(X).
\end{equation}
The composition is multiplication by $M$; it is not asserted
to be the identity integrally.

\emph{Apply this factorization to the family.}
The pullback $(\operatorname{id}_A\times f_j)^*Z$ is a
zero-cycle family on $T_j$. If $a({}^tZ)=0$, its cylinder map
is zero too, by functoriality: it is the original cylinder
map followed by
$f_j^*:H^{2n-3}(X,\mathbb Q)\to H^{2n-3}(T_j,\mathbb Q)$.
The corresponding homomorphism
$A\to\operatorname{Alb}(T_j)$ is therefore zero by faithfulness
on $H_1$, or by the zero-cycle endpoint case E.3d.
Equation \eqref{eq:E3-Alb-factorization} gives $M u_Z=0$.
Finally $A(\mathbb C)$ is divisible: for every $a$ there is
$b$ with $Mb=a$. Hence $u_Z(a)=M u_Z(b)=0$. This proves
the integral family-map converse under the stated support
hypothesis.

\paragraph{E.3k (Refuted: the unqualified integral biconditional in Lead E4;
proved corrected fixed-target criterion).}
For the fourfold $X=W\times E$ of E.3h, the map
$\mathrm{CH}_0(W)\to\mathbb Z$ is an isomorphism by
\citet[Lemma 2.4, author p.~4]{cyOttemSuzuki2020}.
Fix $w_0\in W(\mathbb C)$. For every point $(w,e)$,
the equality $[w]=[w_0]$ on $W$, pushed into $W\times\{e\}$,
shows that its zero-cycle class is represented on
$\{w_0\}\times E$. Thus this curve supports
$\mathrm{CH}_0(X)$, and E.3j proves the family-map converse
for every codimension-three abelian family on this $X$.
E.3h simultaneously gives an algebraically trivial cycle from
an abelian family that is Abel--Jacobi trivial but is not
Albanese trivial. The two assertions declared equivalent in
the literal integral Lead E4 therefore have different truth
values on the same smooth projective fourfold. This refutes
that biconditional, while leaving E.3c's rational formulation
intact.

For an arbitrary smooth projective complex $X$, a precise
fixed-target criterion uses
\[
 K_W^i(X)=\ker\bigl(P^i(X)\longrightarrow J_W^i(X)(\mathbb C)\bigr).
\]
The integral family-map converse is exactly the assertion
that no nonzero image $u_Z(A(\mathbb C))$ of an abelian family
is contained in $K_W^i(X)$. Indeed such containment is
equivalent to $\phi_{W,Z}=0$, which is equivalent to
$a({}^tZ)=0$: one direction uses the finite-kernel argument
in E.3i, and the other the cylinder map in E.3b.
The rational criterion in E.3c is equivalently
$K_W^i(X)\otimes\mathbb Q=0$, because
\[
 0\longrightarrow K_W^i(X)\longrightarrow
 \ker\overline{\operatorname{AJ}}
 \longrightarrow\ker\alpha\longrightarrow0
\]
is exact and $\ker\alpha$ is finite; surjectivity on the right
follows from surjectivity of $\psi_W$.
Neither condition follows merely from regularity or from the
existence of a universal regular homomorphism. We do not
identify $P^i(X)$ with the Walker Jacobian in general.

\paragraph{E.3l (Proved: the integral Walker quotient for one-cycles
under support on a surface).}
Let $X/\mathbb C$ be a smooth projective connected $n$-fold,
$n\ge2$, and assume that $\mathrm{CH}_0(X)$ is supported on
a closed subvariety of dimension at most two. Then the
Walker map induces a canonical isomorphism of abelian groups
\begin{equation}\label{eq:E3-Walker-quotient}
 P^{n-1}(X)\ \xrightarrow{\ \sim\ }\ 
               J_W^{n-1}(X)(\mathbb C).
\end{equation}
In particular the integral family-map converse holds in this
codimension, and its discrepancy from ordinary Abel--Jacobi
equivalence is exactly the finite kernel of
$\alpha:J_W^{n-1}(X)\to J_a^{n-1}(X)$.

\emph{Proof.}
The generic-point, norm, and localization proof in E.3j
applies verbatim with a surface $S$ in place of $C$.
Resolve its two-dimensional components, and lift the
correspondence supported on $X\times S$ by the same
multisection argument. A lift to $X\times\widetilde S$ has
codimension two, so on codimension-$n-1$ algebraically trivial
cycles it acts through
\[
 \mathrm{CH}^1(\widetilde S)_{\mathrm{alg}}
              =\operatorname{Pic}^0(\widetilde S)(\mathbb C).
\]
Lower-dimensional components act through
$\mathrm{CH}^0$ of curves or zero groups and therefore
annihilate algebraically trivial cycles. The term supported
on $D\times X$ still factors through the Albanese varieties
of the resolutions $T_j$ of $D$.

More explicitly, write $j_a:\widetilde S_a\to X$ for the
resolved surface maps. After clearing denominators we can
write the surviving terms of the diagonal as
$(\operatorname{id}_X\times j_a)_*H_a$ and
$(f_j\times\operatorname{id}_X)_*G_j$, where
\[
 H_a\in\mathrm{CH}^2(X\times\widetilde S_a),\qquad
 G_j\in\mathrm{CH}^{n-1}(T_j\times X).
\]
The projection formula, or composition with the respective
graph correspondences, identifies their actions exactly as
$(j_a)_*(H_a)_*$ and $(G_j)_*f_j^*$. Here $f_j^*$ is the
pullback in the Chow ring between smooth varieties, defined
by the graph's refined Gysin map; it is defined even though
$D$ need not be smooth. In particular
\[
\begin{aligned}
 (H_a)_*&: \mathrm{CH}^{n-1}(X)_{\mathrm{alg}}
                \longrightarrow\operatorname{Pic}^0(\widetilde S_a),\\
 f_j^*&: \mathrm{CH}^{n-1}(X)_{\mathrm{alg}}
                \longrightarrow\mathrm{CH}_0(T_j)^0.
\end{aligned}
\]
Composing the second map with $\operatorname{alb}_{T_j}$
therefore gives the asserted factors. Every correspondence
preserves $N$: compose it with each generator's correspondence
in the definition of $N$, leaving the same Albanese-kernel
zero-cycle on its parameter variety. At the two endpoints,
$P^1(\widetilde S_a)=\operatorname{Pic}^0(\widetilde S_a)$ and
$P^{n-1}(T_j)=\operatorname{Alb}(T_j)(\mathbb C)$ by E.3d.
Thus both forward and return maps below are well-defined on
the indicated quotients. After clearing denominators,
we have a factorization of multiplication by an integer $M>0$,
\[
 B=\prod_a\operatorname{Pic}^0(\widetilde S_a)
       \times\prod_j\operatorname{Alb}(T_j),
 \qquad
 P^{n-1}(X)\xrightarrow{F}B(\mathbb C)
 \xrightarrow{G}P^{n-1}(X),
\]
with $GF=M\operatorname{id}$.
Choose an algebraically trivial representative $z$ of a class
in $\ker\overline{\operatorname{AJ}}$. Ordinary Abel--Jacobi
functoriality for correspondences gives, integrally,
\[
\begin{aligned}
 \operatorname{AJ}^1_{\widetilde S_a}((H_a)_*z)
       &=(H_a)_{\operatorname{AJ}}\operatorname{AJ}^{n-1}_X(z)=0,\\
 \operatorname{alb}_{T_j}(f_j^*z)
       &=(f_j^*)_{\operatorname{AJ}}\operatorname{AJ}^{n-1}_X(z)=0.
\end{aligned}
\]
The first Abel--Jacobi map is the identification of
algebraically trivial divisors with $\operatorname{Pic}^0$;
the second is the zero-cycle Abel--Jacobi map, namely the
Albanese map. Hence every component of $F(z)$ is zero.
This uses ordinary Abel--Jacobi functoriality, and does not
require an integral inverse to an isogeny. Consequently
\begin{equation}\label{eq:E3-bounded-AJ-kernel}
             M\ker\overline{\operatorname{AJ}}=0.
\end{equation}

On the other hand, the exact primary statement
\citet[Theorem 2.3, author p.~12]{cyACMVWalker2022} says
that the \emph{whole} kernel of $\psi_W$ is divisible.
Since $N^{n-1}(X)\subset\ker\psi_W$, the group
\[
 K_W^{n-1}(X)=\ker\psi_W/N^{n-1}(X)
\]
is divisible. It is a subgroup of
$\ker\overline{\operatorname{AJ}}$, so
\eqref{eq:E3-bounded-AJ-kernel} kills it by $M$.
A divisible group killed by a positive integer is zero:
write each element as $M$ times another element. Thus
$K_W^{n-1}(X)=0$. Surjectivity of $\psi_W$ proves
\eqref{eq:E3-Walker-quotient} integrally. The exact sequence
in E.3k now identifies the ordinary Abel--Jacobi kernel on
$P^{n-1}(X)$ with $\ker\alpha$ itself.

Under this isomorphism, a family map $u_Z$ is precisely the
regular homomorphism $\phi_{W,Z}$ on complex points. If its
cohomological action is zero, E.3i proves
$\phi_{W,Z}=0$, so $u_Z=0$. This also proves the converse
for the full support-on-a-surface class. It does not extend
the quotient identification to intermediate codimensions
without the stated support hypothesis.

\paragraph{E.3m (Proved: the exact integral Walker reformulation of the
general family converse).}
Fix a smooth projective complex variety $X$ and a codimension $i$.
The following assertions are equivalent:
\begin{enumerate}
\item For every abelian variety $A$ and every normalized
      $Z\in\mathrm{CH}^i(A\times X)$, the vanishing
      $a({}^tZ)=0$ implies $u_Z=0$ integrally.
\item The natural surjection
      $P^i(X)\to J_W^i(X)(\mathbb C)$ is an isomorphism.
\item $N^i(X)=\ker\psi_W$ in $\mathrm{CH}^i(X)_{\mathrm{alg}}$.
\end{enumerate}
These equivalent fixed-target conditions are refuted
universally by E.3u below. Their exact common condition is
the Walker-kernel equality, which differs from the false
integral ordinary Abel--Jacobi equality in Lead E4.

We write $J=J_W^i(X)$, $P=P^i(X)$, and
$\overline\psi:P\twoheadrightarrow J(\mathbb C)$.
The equivalence of (2) and (3) follows from surjectivity of
$\psi_W$ and E.3g's inclusion $N^i\subset\ker\psi_W$.
Condition (2) implies (1) by E.3i. We prove the other direction.

\emph{Choose one family covering the Walker Jacobian.}
Assume (1). Every algebraically trivial cycle is an integral
point of an abelian family by E.3g. The corresponding regular
Walker maps have abelian-subvariety images which generate $J$.
A finite sum of these images is all of $J$: otherwise one can
successively enlarge their dimensions, which is impossible
after $\dim J$ strict enlargements. Thus a product of finitely
many such parameter varieties gives an abelian family
\[
 u_B:B(\mathbb C)\longrightarrow P,\qquad
 \phi_B=\overline\psi\,u_B:B\twoheadrightarrow J.
\]
If $J=0$, (1) immediately kills every family and hence $P$;
we may therefore assume $\dim J>0$.

Let $B_0=(\ker\phi_B)^0$. The restricted family on $B_0$
has zero Walker map and therefore zero cohomological action
by E.3b--i. Condition (1) gives $u_B(B_0(\mathbb C))=0$.
Consequently $u_B$ factors through $D=B/B_0$, and
$D\to J$ is an isogeny. Its abstract kernel on $P$ is a
subgroup of the finite group $\ker(D\to J)(\mathbb C)$.
Over $\mathbb C$ this subgroup is the point group of a finite
subgroup scheme $H\subset D$. Put $E=D/H$. We obtain
\[
 u_E:E(\mathbb C)\hookrightarrow P,
 \qquad \lambda:E\twoheadrightarrow J,
 \qquad \overline\psi\,u_E=\lambda,
\]
where $\lambda$ is an isogeny and $u_E$ is injective. At this
point $u_E$ is a homomorphism of groups; we have not assumed
that a universal integral cycle exists on $E\times X$.

\emph{Every abelian-family image is already this one.}
For any other abelian family $u_A:A(\mathbb C)\to P$, take
the sum family on $B\times A$ and repeat the preceding
construction. Its Walker map is still surjective. We get an
abelian variety $E'$ isogenous to $J$ and an injection
$u_{E'}:E'(\mathbb C)\hookrightarrow P$ whose image is
$u_B(B(\mathbb C))+u_A(A(\mathbb C))$.
The inclusion $B\to B\times A$ induces an algebraic
homomorphism $B\to E'$. It kills the kernel of $B\to E$,
because its composition with the injective map $u_{E'}$
is $u_B$. Quotient descent therefore gives an algebraic
homomorphism $E\to E'$. It is injective on complex points,
again by the injections into $P$. Both abelian varieties have
dimension $\dim J$, so this homomorphism is an isomorphism.
It follows that $u_A(A(\mathbb C))\subset u_E(E(\mathbb C))$.
Using E.3g once more gives
\[
                         P=u_E(E(\mathbb C)).
\]
Moreover, the preceding quotient construction shows that
every composite $u_E^{-1}u_A:A\to E$ is an algebraic
homomorphism, not merely an abstract homomorphism on points.

The kernel of $\overline\psi$ is now the finite group
$\ker\lambda(\mathbb C)$. But it is also
$\ker\psi_W/N^i(X)$, a divisible group by
\citet[Theorem 2.3, author p.~12]{cyACMVWalker2022}.
A finite divisible group is zero. Hence $\lambda$ is an
isomorphism and (2) follows. This already completes the
equivalence; no universality assumption on the Walker map
has entered the argument.

\emph{Regularity of the resulting quotient, also checked directly.}
Let
\[
 q:\mathrm{CH}^i(X)_{\mathrm{alg}}\longrightarrow P
                       \xrightarrow{u_E^{-1}}E(\mathbb C).
\]
For a pointed smooth projective parameter variety $(T,t_0)$
and a correspondence $W$, its point-to-cycle map modulo $N$
factors integrally through $\operatorname{Alb}(T)(\mathbb C)$,
because $W_*$ kills the Albanese kernel by the definition of
$N$. Denote the resulting homomorphism to $E(\mathbb C)$ by
$v$. E.3a constructs an integer $m>0$ and an abelian family
on $\operatorname{Alb}(T)$ whose map to $P$ is $m$ times
the original one. Thus $mv$ is an algebraic homomorphism
$f:\operatorname{Alb}(T)\to E$.

We check that $v$ itself is algebraic. For every
$a\in\operatorname{Alb}(T)[m](\mathbb C)$,
$f(a)=mv(a)=v(ma)=0$. In characteristic zero this means
that $f$ kills the whole finite group scheme
$\operatorname{Alb}(T)[m]$. Hence $f$ factors algebraically
through $[m]$ as $f=g\circ[m]$ for a homomorphism
$g:\operatorname{Alb}(T)\to E$. The difference $v-g$ takes
values in $E[m](\mathbb C)$; its source is divisible, so this
finite-valued homomorphism is zero. Therefore $v=g$.
Composing with the pointed Albanese map proves regularity
for projective parameter varieties.

For a smooth quasi-projective parameter variety, choose a
smooth projective compactification and close the correspondence;
the refined fibres on the original open are unchanged. This
reduces regularity to the preceding case. General smooth
finite-type parameter schemes can be covered by such opens,
and the resulting morphisms agree because their complex-point
maps agree and the target is separated. Thus $q$ is a regular
homomorphism in the usual sense.

We have, in particular,
\[
                    \psi_W=\lambda\circ q.
\]
The direct regularity check is consistent with the alternative
isogeny-initiality proof from the same Theorem 2.3, but that
alternative is unnecessary after the finite-divisible-kernel
argument. The assertion that the original family converse
holds for every $X,i$ remains unproved here.

\paragraph{E.3n (Proved: a torsion-valued abelian-family map to $P$ over
$\mathbb C$ is zero).}
Let $X/\mathbb C$ be smooth projective, let $A$ be an abelian
variety, and let $Z\in\mathrm{CH}^i(A\times X)$ be normalized.
Then
\[
 u_Z\otimes\mathbb Q=0\quad\Longrightarrow\quad u_Z=0.
\]
The hypothesis means that every value of $u_Z$ is torsion;
it does not assume a uniform exponent. We prove the required
uniformity from finite-type parameter spaces for the defining
Albanese relations. Uncountability of $\mathbb C$ is essential
to this proof, as the example E.3f over $\overline{\mathbb F}_p$
also shows.

\emph{Moving the points before parameterizing the relations.}
For a smooth projective connected $T$ and a nonempty open
$T^\circ\subset T$, every class in $\mathrm{CH}_0(T)$ has a
zero-cycle representative supported in $T^\circ$; equivalently,
the map $Z_0(T^\circ)\to\mathrm{CH}_0(T)$ from the group of
cycles is surjective. We are not using a covariant Chow-group
map for an open immersion.
Indeed, for each excluded point choose an integral curve
through it whose generic point lies in $T^\circ$, normalize
the curve, and use Riemann--Roch to express that point as a
difference of effective divisors avoiding the finitely many
points outside $T^\circ$. Pushing to $T$ proves the assertion.
For dimension zero it is immediate.

Represent a correspondence $W\in\mathrm{CH}^i(T\times X)$
by a finite signed sum of integral subvarieties. Those which
do not dominate $T$ act trivially on $\mathrm{CH}_0(T)$:
move the zero-cycle away from their images. Each dominating
component is flat over a common nonempty open $T^\circ$
after generic flatness. We can therefore replace any
Albanese-kernel zero-cycle $\nu$ by a rationally equivalent
zero-cycle supported in $T^\circ$. Its Albanese class remains
zero. Now $W_*\nu$ is represented by the actual signed
sum of flat fibres $W_t$ at its finitely many support points.
Thus every generator of $N^i(X)$ has this form; we need not
parameterize arbitrary refined intersections at bad fibres.

\emph{Countably many finite-type spaces parameterize these generators.}
Embed $X$ once in projective space. We describe all data as
the projective embedding dimension, Hilbert polynomials,
finite lengths, and integer coefficients vary. There are
only countably many such choices.
Throughout we replace the parameter bases by their reductions
and use finite locally closed stratifications when required
for Hilbert--Chow, constant fibre dimension, or multiplicity.
This does not lose any complex-point datum; allowing all the
strata preserves countability and constructibility.

First use Hilbert schemes to parameterize closed embedded
smooth projective geometrically connected varieties $T$.
The locus with these fibre properties is open in the
appropriate Hilbert scheme; allowing the finitely many
components and locally closed strata gives finite-type bases.
Pass to the universal $T$ itself as a base to mark an origin
$t_0$. The pulled-back family has a section. Its relative
Albanese is an abelian scheme, with its pointed map, compatible
with base change. The precise characteristic-zero existence
and base-change statement is
\citet[Theorem 2.2]{cyJavanpeykarMathur2022}; the marked section
turns their Albanese torsor into the pointed group scheme.

Relative Hilbert schemes for $T\times X$ parameterize the
finitely many components of $W$. Restrict to the locus where
the marked evaluation points lie in the fibre-flat locus.
This is an algebraic constructible condition. More explicitly,
the universal subschemes are flat over their Hilbert base;
the local criterion for flatness identifies fibrewise flatness
over $T$ with flatness of their projections over the universal
$T$, and the bad locus has closed image because these
projections are projective. Finite stratification also fixes
the dimensions and multiplicities of the resulting fibre
cycles. Mark the finitely many points $t_1,\ldots,t_s$ by
the relative power $T^s$. For fixed integer weights $e_j$
with $\sum_j e_j=0$, the condition
\[
            \sum_j e_j\operatorname{alb}_{T}(t_j)=0
\]
is closed: it is the inverse image of the zero section under
an algebraic morphism to the relative Albanese. On this
finite-type locus, the signed flat-fibre cycle
$\sum_j e_jW_{t_j}$ is an algebraic family of cycles on $X$.
Its positive and negative parts are represented, after finite
stratification, by morphisms to fixed-degree Chow schemes.
Equivalently one can apply the Hilbert--Chow morphism to each
flat family and then use the algebraic addition of cycles.
Allowing finite tuples of these data parameterizes every
finite sum of generators of $N^i(X)$.

\emph{Rational-equivalence witnesses also have countable finite-type data.}
An integral rational equivalence is a finite sum of signed
differences $[V_0]-[V_\infty]$, where
$V\subset\mathbb P^1\times X$ is an integral closed subscheme
dominating $\mathbb P^1$. Such a subscheme is flat over the
smooth curve $\mathbb P^1$. The equivalence of this description
with principal-divisor rational equivalence follows by closing
the graphs of the finitely many rational functions and pushing
forward; it is proved in
\citet[Lemma 43.9.1, Tag 0AZI]{cyStacksRationalEquivalence}.
Hilbert schemes of $\mathbb P^1\times X$, indexed by their
countably many Hilbert polynomials, therefore parameterize
all witnesses. We restrict to the locus of flat projection
to $\mathbb P^1$, evaluate at $0,\infty$, and apply
Hilbert--Chow and cycle addition. Equality of the two resulting
effective sums is a closed condition in a fixed Chow scheme.
Finite tuple lengths and integer signs give only countably
many more choices. Consequently the locus that two
specified algebraic cycle families are rationally equivalent
is a countable union of constructible sets, by Chevalley's
theorem. Applying this to the preceding finite sums proves
that membership in $N^i(X)$, in any fixed algebraic family
of signed cycles, is countably constructible.

For the same Hilbert--Chow construction of countable
rational-equivalence witnesses, with its maps written out,
see the opened primary proof of
\citet[Lemma 4.4, pp.~311--313]{cySchoen1993}.
Our additional parameter spaces above impose the relative
Albanese-kernel condition; that condition is not asserted
to be part of Schoen's lemma.

All varieties and cycles in a particular relation occur in
this enumeration: projective embeddings and Hilbert polynomials
are finite data, the Albanese condition is exact, the moving
step preserves the given Chow class integrally, and every
rational equivalence uses finitely many functions. Thus the
construction does not replace $N$ by its saturation, closure,
or a subgroup defined only by regular invariants.

\emph{Apply uncountability and the group law.}
Choose a nonempty Zariski open $U\subset A$ where the finitely
many dominating components of a representative of $Z$ are
flat; vertical components vanish on the fibres there. Fix
any cycle representative of its refined fibre at the origin.
Then $Z_a-Z_0$ is a signed algebraic cycle family for $a\in U$.
For every $m>0$, the preceding construction gives
\[
 \{a\in U(\mathbb C):m u_Z(a)=0\}
                 =\bigcup_{j\ge1}C_{m,j}(\mathbb C),
\]
where every $C_{m,j}$ is constructible in $U$.
If every $u_Z(a)$ is torsion, the countable collection of all
$C_{m,j}$ covers $U(\mathbb C)$. A complex irreducible variety
cannot be the union of countably many proper closed algebraic
subsets: this follows, for example, from the Baire theorem on
a coordinate ball. Therefore some $C_{m,j}$ is Zariski dense
and contains a nonempty open $V\subset A$.

The full kernel of $m u_Z$ is a subgroup of $A(\mathbb C)$,
so it contains $V(\mathbb C)-V(\mathbb C)=A(\mathbb C)$.
The equality follows because $V$ meets every translate of
$V$ in the irreducible variety $A$. Thus $m u_Z=0$ with one
fixed positive integer $m$. Divisibility of $A(\mathbb C)$
then gives $u_Z=0$, as claimed.

Combining this result with E.3c and E.3m yields an exact
further reduction: for fixed $X,i$, the integral family
converse, the rational family converse, the rational
ordinary Abel--Jacobi injectivity on $P^i(X)\otimes\mathbb Q$,
and the integral equality $N^i(X)=\ker\psi_W$ are equivalent.
Their universal truth is refuted by the specified
codimension-four family of E.3u. Ordinary integral Abel--Jacobi
injectivity is a different assertion, refuted in E.3h--k.

\paragraph{E.3o (Known, with the integral consequence proved: the Walker
kernel is Saito's product relation).}
We work over $\mathbb C$ with integral Chow groups. For a smooth
projective variety $V$, let
\[
 \mathcal A^r(V)=\mathrm{CH}^r(V)_{\mathrm{alg}},\qquad
 \mathcal H^r(V)=\ker\bigl(\mathrm{CH}^r(V)
                 \longrightarrow H^{2r}(V,\mathbb Z(r))\bigr).
\]
For two adequate relations $\mathcal E,\mathcal F$, we use
Saito's product: $(\mathcal E*\mathcal F)^i(X)$ is the subgroup
generated by
\[
 (p_X)_*(a\cdot b),\qquad
 a\in\mathcal E^r(T\times X),\quad
 b\in\mathcal F^s(T\times X),\quad
 r+s=\dim T+i,
\]
where $T$ is smooth projective and $p_X:T\times X\to X$.
This is an adequate relation by
\citet[Section 1.3 and Lemma 1.3.1, pp.~293--294]{cySaito1992}.
For an adequate relation $\mathcal E$, we also set
\[
 \langle\mathcal E\rangle_0^i(X)
 =\sum_{T,\,W\in\mathrm{CH}^i(T\times X)}
       W_*\bigl(\mathcal E^{\dim T}(T)\bigr).
\]
This is exactly the zero-cycle-generated relation of
\citet[Section 1.7 and Lemma 1.7.1, p.~296]{cySaito1992}.
The following equalities are integral:
\begin{equation}\label{eq:walker-saito-product}
 \ker\psi_W^i(X)=(\mathcal A*\mathcal H)^i(X),
 \qquad
 N^i(X)=\langle\mathcal H^{*2}\rangle_0^i(X).
\end{equation}
Consequently the general complex converse in E.3 is equivalent
to the precise generation assertion
\begin{equation}\label{eq:remaining-zero-generation}
 \boxed{\quad
 \langle\mathcal H^{*2}\rangle_0^i(X)
       =(\mathcal A*\mathcal H)^i(X).
 \quad}
\end{equation}
E.3u below refutes \eqref{eq:remaining-zero-generation}
in general, by an explicit codimension-four family.

\begin{proof}
We first identify exactly the statement we take from Saito.
His Example 1.8, equation (1.8.2) on p.~296, defines
\[
 \operatorname{Gr}^{l}\mathrm{CH}(V)
 =\frac{\langle\mathcal H^{*l}\rangle_0\mathrm{CH}(V)}
 {\bigl(\langle\mathcal H^{*l}\rangle_0
                 \cap\mathcal H^{*(l+1)}\bigr)\mathrm{CH}(V)}.
\]
In particular its first step is
$\mathcal A/(\mathcal A\cap\mathcal H^{*2})$, since
$\langle\mathcal H\rangle_0=\mathcal A$ by his Example 1.9,
p.~297. The denominator is an intersection; it is not
$\langle\mathcal H^{*(l+1)}\rangle_0$.
Thus Theorem 6.2 does not by itself identify the manuscript's
quotient $P^i(X)$ with an abelian variety.

The needed positive statement is explicit in
\citet[proof of Theorem 6.2, p.~323; Remark 6.4, p.~325]{cySaito1992}:
\[
 K_{\mathrm{AJ}}^i(X)/(\mathcal A*\mathcal H)^i(X)
 \quad\text{is finite},\qquad
 K_{\mathrm{AJ}}^i(X)=
       \ker\bigl(\mathcal A^i(X)\to J_a^i(X)(\mathbb C)\bigr).
\]
For clarity, Saito introduces
$'\operatorname{Gr}^{1}\mathrm{CH}^i(X)
=\mathcal A^i(X)/(\mathcal A*\mathcal H)^i(X)$ in
Lemma 6.2.1, p.~321. Its map to $J_a^i(X)$ is covered by
an abelian variety whose cohomological map is an isomorphism
over $\mathbb Q$; the composition is an isogeny. The
surjectivity proved on pp.~322--323 makes its kernel finite.
This is an assertion about the integral groups just displayed,
not a statement only after tensoring by $\mathbb Q$.

We record why the two divisibility steps have the required
scope. The group $(\mathcal A*\mathcal H)^i(X)$ is divisible:
in each generator $(p_X)_*(a b)$ we divide the algebraically
trivial factor $a$ by the chosen positive integer, leaving
$b$ fixed. This also proves directly the instance of Saito's
Lemma 1.10, p.~297, that we use. Ordinary Abel--Jacobi
vanishes on this group. Equivalently this follows from the
inclusion in Remark 6.4, or by representing $a$ by a curve
family and using that multiplication by $b$ makes the
correspondence's cohomological action zero.

Write
\[
 \mathcal A^i(X)\xrightarrow{\psi_W^i}J_W^i(X)(\mathbb C)
     \xrightarrow{\alpha}J_a^i(X)(\mathbb C),
\]
where $\alpha$ is the Walker isogeny. The image of
$(\mathcal A*\mathcal H)^i(X)$ under $\psi_W^i$ lies in
the finite group $\ker\alpha(\mathbb C)$. A homomorphism
from a divisible group to a finite group is zero; hence
$(\mathcal A*\mathcal H)^i(X)\subset\ker\psi_W^i$.
The quotient
\[
 \ker\psi_W^i/(\mathcal A*\mathcal H)^i(X)
\]
is therefore a subgroup of Saito's finite quotient. It is
also divisible, because the whole Walker kernel is divisible
by \citet[Theorem 2.3, author p.~12]{cyACMVWalker2022}.
It is zero. This proves the first equality in
\eqref{eq:walker-saito-product}.

For a smooth projective $t$-fold $T$, Saito's Corollary 6.3
and Remark 6.4, p.~325, specialize to the integral equality
\[
 (\mathcal A*\mathcal H)^t(T)
 =\ker\bigl(\mathrm{CH}_0(T)^0
                      \to\operatorname{Alb}(T)(\mathbb C)\bigr).
\]
Here the absence of a finite discrepancy uses the injectivity
of the zero-cycle Albanese map on torsion, explicitly used
in his proof on p.~323. Moreover
$\mathcal A*\mathcal H\subset\mathcal H^{*2}$, and
$\mathcal H^{*2}$ is Abel--Jacobi trivial by Example 1.11,
pp.~297--298. On zero-cycles, ordinary Abel--Jacobi is the
Albanese map. These two inclusions give
\[
 (\mathcal H^{*2})^t(T)
 =\ker\bigl(\mathrm{CH}_0(T)^0
                      \to\operatorname{Alb}(T)(\mathbb C)\bigr).
\]
Taking the correspondence images in the definition of
$\langle\mathcal H^{*2}\rangle_0$ is now precisely the
definition of $N^i(X)$. This proves the second equality.
Together with E.3m--n the two equalities prove the asserted
fixed-target equivalence.
\end{proof}

The same product $\mathcal A*\mathcal H$ is obtained if
the second factor is required only to have zero rational
Betti class. Indeed an integral cycle $b$ with zero rational
class has $mb\in\mathcal H$ for some $m>0$; writing
$a=ma'$ within the divisible group $\mathcal A$ gives
$ab=a'(mb)$. This observation does not replace the
integral homological relation in the displayed definition
of $\langle\mathcal H^{*2}\rangle_0$.

\paragraph{E.3p (Proved: only surfaces are needed in the
generation test).}
For every smooth projective complex $X$ and every $i$,
\begin{equation}\label{eq:surface-generated-N}
 N^i(X)=
 \sum_{\substack{S/\mathbb C\ \mathrm{smooth\ projective\ surface}\\
                 W\in\mathrm{CH}^i(S\times X)}}
 W_*\ker\bigl(\mathrm{CH}_0(S)^0
                   \to\operatorname{Alb}(S)(\mathbb C)\bigr)
 \quad\text{integrally}.
\end{equation}
In particular \eqref{eq:remaining-zero-generation} is the
assertion that every sum of the products in E.3o is a finite
sum of these surface-kernel images. This is a restriction
on the dimension of the generating varieties, not an
assertion that their Albanese kernels are themselves known.
It agrees with
\citet[Lemma 7.1.1, pp.~326--327]{cySaito1992}; we give
the direct argument appropriate to the present definition.

\begin{proof}
Consider one generator $W_*\nu$ of $N^i(X)$, where
$T$ is smooth projective connected and $\nu\in
\mathrm{CH}_0(T)^0$ has zero Albanese class. Choose an
integral zero-cycle representative of $\nu$ and let $F$
be its finite support. If $\dim T\ge3$, sufficiently
positive successive hyperplane sections through $F$
give a smooth connected complete-intersection surface
$j:S\hookrightarrow T$. To see that the imposed points
cause no obstruction, take powers of a very ample bundle
whose sections vanishing on $F$ generate the conormal
directions at every point of $F$; a general section is
then smooth there, and Bertini applies away from $F$.
Repeat until the dimension is two.

Integral weak Lefschetz gives
$H_1(S,\mathbb Z)\simeq H_1(T,\mathbb Z)$, compatibly
with the Hodge structures and their periods. Thus
\[
 j_*:\operatorname{Alb}(S)\xrightarrow{\sim}
                         \operatorname{Alb}(T)
\]
is an isomorphism of complex abelian varieties, rather
than merely an isogeny. The same formal sum of the points
of $F$ defines a degree-zero cycle $\nu_S$ on $S$.
Its Albanese class maps to that of $\nu$, so is zero.
The refined pullback $W_S=(j\times1_X)^*W$ is an
integral class in $\mathrm{CH}^i(S\times X)$, and
the compatibility of correspondences with proper
pushforward gives
\[
                   W_*\nu=(W_S)_*\nu_S.
\]
No transversality of a chosen representative of $W$
is required for this equality of Chow operations.

If $\dim T=2$ there is nothing to do. In dimension one
the Albanese kernel is zero by Abel's theorem, and
in dimension zero the connected degree-zero Chow group
is zero. We have expressed every generator using a
surface; the reverse inclusion is immediate from the
definition of $N$.
\end{proof}

The literal ordinary Abel--Jacobi equality is still
false in the Ottem--Suzuki example of E.3h--k. The
zero-cycle generation assertion in
\eqref{eq:remaining-zero-generation}, equivalently the
surface generation in \eqref{eq:surface-generated-N}
of all $\mathcal A*\mathcal H$ classes, is refuted by E.3u.
A universal
regular-homomorphism theorem does not assert this
equality of kernels: Samuel's abelian equivalence is
defined by the intersection of all regular kernels,
whereas $N$ is generated by the specified correspondences.
The distinction is explicit in
\citet[Section 1.1, pp.~7--8]{cyBarbieriViale2001}
and \citet[Section 4.1, pp.~13--16]{cyHernandezPelaez2025}.
\paragraph{E.3q (Proved under stated hypotheses: the codimension-two
case with surface-supported zero-cycles).}
Let $X/\mathbb C$ be smooth projective of dimension $n\ge2$,
and suppose that $\mathrm{CH}_0(X)$ is supported on a closed
subset of dimension at most two, in the sense used in E.3l.
Then
\[
 P^2(X)\ \xrightarrow{\ \sim\ }\ J_W^2(X)(\mathbb C)
          \ \xrightarrow{\ \sim\ }\ J_a^2(X)(\mathbb C)
\]
integrally. Thus the family converse and the integral equality
of ordinary Abel--Jacobi and Albanese equivalence hold in
codimension two under this hypothesis. Together with E.3l this
settles both codimensions $2$ and $n-1$ under the same support
assumption; for $n=3$ they are the same codimension.

\begin{proof}
We transpose the actual decomposition of the diagonal
constructed in E.3l. After resolving its support and clearing
one integer, its action on $\mathrm{CH}^2(X)$ has the form
\begin{equation}\label{eq:transpose-surface-factorization}
 M\,\mathrm{id}
   =\sum_j (G_j)_* f_j^*
          +\sum_k (g_k)_* (H_k)_* .
\end{equation}
Here $M>0$, the smooth projective $S_j$ have dimension
$s_j\le2$, and the smooth projective $D_k$ have dimension
$n-1$. The proper maps and correspondences have the types
\[
 \begin{gathered}
 f_j:S_j\longrightarrow X,\qquad
 G_j\in\mathrm{CH}^{s_j}(S_j\times X),\\
 g_k:D_k\longrightarrow X,\qquad
 H_k\in\mathrm{CH}^{n-1}(X\times D_k).
 \end{gathered}
\]
Consequently the first summand in
\eqref{eq:transpose-surface-factorization} acts through
$\mathrm{CH}^2(S_j)$, and the second acts through
$\mathrm{CH}^1(D_k)$:
\[
 \begin{gathered}
 \mathrm{CH}^2(X)\xrightarrow{f_j^*}
       \mathrm{CH}^2(S_j)\xrightarrow{(G_j)_*}\mathrm{CH}^2(X),\\
 \mathrm{CH}^2(X)\xrightarrow{(H_k)_*}
       \mathrm{CH}^1(D_k)\xrightarrow{(g_k)_*}\mathrm{CH}^2(X).
 \end{gathered}
\]
All these operations are integral and preserve algebraic
equivalence and $N$. The same support-resolution argument
as in E.3l supplies the integer $M$; transposition does not
introduce any additional denominator.

Let $z\in\mathcal A^2(X)$ have zero ordinary Abel--Jacobi
class. Correspondence functoriality implies that
$f_j^*z$ has zero Albanese class when $s_j=2$.
It is zero for dimension reasons when $s_j<2$.
Likewise $(H_k)_*z\in\operatorname{Pic}^0(D_k)$ has
zero divisor Abel--Jacobi class, so is zero as a
divisor class. Formula \eqref{eq:transpose-surface-factorization}
therefore gives $Mz\in N^2(X)$. It follows that
$\ker\psi_W^2/N^2(X)$ is killed by $M$.
It is a divisible group by ACMV's Theorem 2.3, and hence
zero. This proves the first isomorphism.

For the second, Saito's Corollary 6.3 and Remark 6.4,
p.~325, give
$\ker\psi_{\mathrm{AJ}}^2=(\mathcal A*\mathcal H)^2(X)$
integrally; E.3o identifies this with $\ker\psi_W^2$.
Since $\psi_W^2$ is surjective, the Walker isogeny has
zero kernel and is an isomorphism. This argument uses
the exact codimension-two torsion theorem incorporated
in Saito's Corollary 6.3; it does not extend the same
ordinary-Jacobian conclusion to codimension $n-1$ in
higher dimensions.
\end{proof}



\paragraph{E.3r (Proved: a specified generic fourfold and a two-torsion
counterexample).}
We can remove the unspecified auxiliary prime and generality choices
from the variety in E.3h--k. Put
\[
 \mathcal M_2=\{\alpha\in\mathbb Z_{\ge0}^3:|\alpha|=2\},
 \qquad
 k=\overline{\mathbb Q(a_{j\epsilon\alpha},b_{j\epsilon\alpha})},
\]
where all $72$ elements $a_{j\epsilon\alpha},b_{j\epsilon\alpha}$,
with $1\le j\le3$, $0\le\epsilon\le1$, and
$\alpha\in\mathcal M_2$, are algebraically independent.
On $\mathbb P^1_k\times\mathbb P^2_k\times\mathbb P^2_k$,
with coordinates $[S:T]$, $[x_0:x_1:x_2]$, $[y_0:y_1:y_2]$, set
\[
 A_j=\sum_{\epsilon,\alpha}
      a_{j\epsilon\alpha}S^{1-\epsilon}T^\epsilon x^\alpha,
 \qquad
 B_j=\sum_{\epsilon,\alpha}
      b_{j\epsilon\alpha}S^{1-\epsilon}T^\epsilon y^\alpha.
\]
Define $W/k$ by the three $2$-by-$2$ minors of
\[
                  \begin{pmatrix}A_1&A_2&A_3\\B_1&B_2&B_3\end{pmatrix}.
\]
Take $t$ transcendental over $k$, put $K=\overline{k(t)}$, and define
\[
 E_t:\ v^2=u(u-1)(u-t),\qquad o=\infty,\qquad e=(0,0),
 \qquad X=W_K\times_K E_t.
\]
For every embedding $K\hookrightarrow\mathbb C$, $X_{\mathbb C}$
is a smooth projective fourfold with $\mathrm{CH}_0$ supported on
$\{w_0\}\times E_t$, for any geometric point $w_0$ of $W$.
There is an integral cycle $\gamma\in\mathrm{CH}^2(W)$ such that
\[
 z=\gamma\times([e]-[o])
       \in\mathrm{CH}^3(X_{\mathbb C})_{\mathrm{alg}}
\]
satisfies
\begin{equation}\label{eq:explicit-os-torsion}
             2z=0,\qquad z\ne0,\qquad
             \operatorname{AJ}(z)=0,\qquad z\notin N^3.
\end{equation}
Below we give an exact instruction for choosing $\gamma$.
The displayed variety has no unselected sufficiently-large integer
in its equations.

\begin{proof}
\emph{Integral transfer of the parity obstruction.}
Let $V=\mathbb A^{72}_{\overline{\mathbb Q}}$ be the coefficient
space. A nonempty Zariski open $U\subset V$ carries a smooth
projective family $\pi:\mathcal W\to U$ of these threefolds.
We choose $U$ so that the first-row zero locus in
$\mathbb P^1\times\mathbb P^2$ is finite etale of degree $12$.
Its product with the last $\mathbb P^2$ is then a relative
effective Cartier divisor $\mathcal D\subset\mathcal W$.
The degree follows from
\[
          \int_{\mathbb P^1\times\mathbb P^2}(h+2l)^3=12.
\]
These are open conditions. The general matrices in
\citet[Lemmas 2.3--2.7 and Theorem 3.1, author pp.~4--8]
{cyOttemSuzuki2020} satisfy them. We take $U$ containing one
of their rational matrices $v_0$. Its divisor $D_{v_0}$ is
their $\sum_{j=1}^{12}E_{1,j}$. Their integral congruence (1),
author p.~7, says
\begin{equation}\label{eq:os-parity-all-curves}
 \deg(c/\mathbb P^1)-D_{v_0}\cdot c=0\pmod2
               \quad(c\in\mathrm{CH}_1(W_{v_0})).
\end{equation}
We use this exact conclusion of their degeneration proof as a
known input. Their auxiliary prime is not part of our
definition of $W$.

Let $C$ be an integral curve on $W_{\mathbb C}$. It defines a
point of a relative Hilbert scheme
$\operatorname{Hilb}_P(\mathcal W/U)$ for the fixed Segre
embedding. This scheme is projective over $U$ and defined over
$\overline{\mathbb Q}$. Choose an irreducible component $H$
containing the point. The coefficient tuple is algebraically
independent, so its image in $U$ is the generic point as an
$\overline{\mathbb Q}$-scheme. Thus $H$ dominates $U$.
Its image is closed by properness, so $H$ has a point over $v_0$.
The flat limit defines a one-cycle $c_0$ with integral multiplicities; embedded or
zero-dimensional components do not affect that one-cycle.
The degrees of both
$\mathcal O_{\mathbb P^1}(1)$ and
$\mathcal O_{\mathcal W}(\mathcal D)$ on a flat
one-dimensional family are constant on each connected Hilbert
component. They can be read as the linear coefficients of
Euler characteristics after tensoring by powers of the line bundle.
Consequently
\[
 \deg(C/\mathbb P^1)-D\cdot C
   =\deg(c_0/\mathbb P^1)-D_{v_0}\cdot c_0\equiv0\pmod2.
\]
Additivity proves this for every integral one-cycle on
$W_{\mathbb C}$. No norm, extension degree, or multiplication
of a cycle occurs in this argument.

The same source, author p.~8, gives
$\beta_0\in H_2(W_{v_0},\mathbb Z)$ with
\[
        \deg(\beta_0/\mathbb P^1)=1,\qquad D_{v_0}\cdot\beta_0=0
\]
after a possible sign change. The base $U(\mathbb C)$ is path
connected. Parallel transport and Poincare duality give
$\beta\in H^4(W_{\mathbb C},\mathbb Z)$ with the same two
intersection numbers. Integral cohomology is torsion-free by
Lemma 2.7 or transport from $v_0$. Lemma 2.3 gives
$H^2(W,\mathcal O_W)=0$, so $H^4(W,\mathbb Z)$ consists of
Hodge classes. The parity congruence shows that $\beta$ is
not algebraic.

\emph{An algebraic lift of a class of order two in the defect.}
The subgroup
\[
 B=H^4_{\mathrm{alg}}(W_{\mathbb C},\mathbb Z)
       \subset H^4(W_{\mathbb C},\mathbb Z)
\]
has finite index: the Lefschetz $(1,1)$ theorem makes
$H^2(W,\mathbb Q)$ algebraic, and multiplication by the ample
class of $\mathcal O(1,1,1)|_W$ is an isomorphism
$H^2(W,\mathbb Q)\to H^4(W,\mathbb Q)$.
Let $m>0$ be least with $m\beta\in B$. Parity makes $m$ even.
Choose any algebraic representative
\[
             \operatorname{cl}(\gamma)=m\beta .
\]
We can take $\gamma$ over $k$: for each component of a complex
representative, choose a $k$-point on its Hilbert component for
the fixed variety $W/k$. Algebraic closedness of $k$ gives
geometric connectedness of the component and existence of such
a point. The resulting algebraically equivalent cycle has the
same integral cohomology class.

The element $\operatorname{cl}(\gamma)\otimes\tfrac12$ in
$B\otimes\mathbb Q/\mathbb Z$ is nonzero. Otherwise
$\operatorname{cl}(\gamma)\in2B$; torsion-freeness of $H^4$
would put $(m/2)\beta$ in $B$, contradicting minimality.
Its image in $H^4(W,\mathbb Z)\otimes\mathbb Q/\mathbb Z$ is
zero because $m\beta$ is even. In particular
$\gamma\otimes\tfrac12$ is nonzero in
$\mathrm{CH}^2(W)\otimes\mathbb Q/\mathbb Z$.

\emph{The specified elliptic curve and cycle.}
The Legendre formulas
\[
 \Delta=16t^2(1-t)^2,\qquad
 j(E_t)=256\,\frac{(1-t+t^2)^3}{t^2(1-t)^2}
\]
show smoothness and that $j(E_t)$ is transcendental over $k$.
The class $\tau=[e]-[o]$ is nonzero of order two in
$\operatorname{Pic}^0(E_t)$. Choose a direct summand
$\mathbb Q/\mathbb Z$ of its torsion group taking $1/2$ to $\tau$.
Then $\gamma\otimes\tau\ne0$.
The exact exterior-product injectivity theorem of
\citet[Theorem (0.2), p.~177]{cySchoen2000Exterior}, for
$k\subset\mathbb C$, gives $z=\gamma\times\tau\ne0$.
This cycle has order two and is algebraically trivial.

The Deligne cohomology exact sequence for $W$ has divisible
kernel $J^2(W)$. After tensoring by
$\operatorname{Pic}^0(E_t)_{\mathrm{tors}}\simeq
(\mathbb Q/\mathbb Z)^2$, this kernel contributes zero.
The Deligne class of the product therefore factors through
\[
 H^4_{\mathrm{Hdg}}(W,\mathbb Z)
          \otimes\operatorname{Pic}^0(E_t)_{\mathrm{tors}}.
\]
Its value is $(m\beta)\otimes\tau=0$.
Thus $\operatorname{AJ}(z)=0$. This also verifies the mechanism
in \citet[Proposition 3.1 and proof, author pp.~5--6]
{cySuzuki2020Remark} for our specified field and elliptic curve.

The same chosen $\gamma$ works under every embedding into
$\mathbb C$. Its nonzero tensor class over $k$ is intrinsic.
Under the embedding used to choose it, its Betti class is even,
so comparison makes its mod-two etale cycle class over $k$ zero.
This is an intrinsic class over the algebraically closed field
$k$. Comparison after any other embedding therefore again makes
its integral Betti class even. Schoen's injectivity and the
Deligne-product argument apply unchanged. This proves the
stated quantifier without choosing a new $\gamma$ for each
embedding.

Finally Ottem--Suzuki's Lemma 2.4 gives
$\mathrm{CH}_0(W_{\mathbb C})=\mathbb Z$. In each
$W\times\{e'\}$ we can move every point to $(w_0,e')$, so
$\{w_0\}\times E_t$ supports $\mathrm{CH}_0(X_{\mathbb C})$.
The hypotheses of
\citet[Theorem 1.1, author p.~2]{cySuzuki2020Remark} hold.
That theorem makes the Walker map injective on algebraically
trivial codimension-three torsion. Hence $\psi_W(z)\ne0$;
E.3g gives $N^3\subset\ker\psi_W$, so $z\notin N^3$.

E.3j applies to this same displayed fourfold and proves its
integral family-map converse. We have therefore specified a
function-field counterexample to the literal integral
biconditional in Lead E4. The rational equivalence in E.3c
remains valid. The exact fixed-target Walker-generation
equivalences of E.3o--p remain valid as well; E.3u separately
refutes their universal conditions.
\end{proof}


% Full source/proof audit passed by root and frobenius.
\paragraph{E.3s (Proved: the exact large one-motivic quotient, without a finiteness assertion).}
We work over $k=\mathbb C$, with integral coefficients, and fix a
smooth projective variety $X$. Let $\mathcal C_X^i$ be the etale
sheaf with transfers associated to the presheaf
\[
                  T\longmapsto\operatorname{CH}^i(T\times X)
\]
on smooth varieties. Write $\operatorname{Alb}_{\mathrm{sh}}$
for the left adjoint to the inclusion of one-motivic etale
sheaves into etale sheaves with transfers, and set
\[
 \mathcal E_X^i=\operatorname{Alb}_{\mathrm{sh}}(\mathcal C_X^i),
 \qquad
 (\mathcal E_X^i)^0=\ker(\mathcal E_X^i\longrightarrow
                                      \pi_0\mathcal E_X^i).
\]
The superscript zero here refers to this sheaf kernel; it does
not presuppose representability by a connected group scheme.
Then the canonical unit has the exact point-group description
\begin{align}
 \ker\bigl(\operatorname{CH}^i(X)\longrightarrow
                           \mathcal E_X^i(\mathbb C)\bigr)
     &=N^i(X),\label{eq:E3-sheaf-kernel}\\
 P^i(X)&\simeq(\mathcal E_X^i)^0(\mathbb C).
                                      \label{eq:E3-sheaf-connected}
\end{align}
In particular this construction gives a sheaf realization of
the manuscript's generated relation, not an identification
with the intersection of all regular kernels.

\begin{proof}
We state the exact categorical inputs. In
\citet[Corollary 1.3.5, Proposition 1.3.11, and
Lemma 1.3.12, author pp.~18, 24--25]{cyAyoubBarbieriViale2009},
the category of one-motivic etale sheaves is a cocomplete
Serre subcategory of etale sheaves with transfers. For a
smooth variety $T$, its Serre--Albanese group scheme
$\operatorname{Alb}(T)$ belongs to that subcategory,
\[
 \mathbb Z_{\mathrm{tr}}(T)\longrightarrow
                           \operatorname{Alb}(T)
\]
is an epimorphism of etale sheaves, and every map from
$\mathbb Z_{\mathrm{tr}}(T)$ to a one-motivic sheaf factors
uniquely through it. Epimorphy is their Lemma 1.3.3,
author pp.~16--17. There is no characteristic exponent to
invert over $\mathbb C$.

We can present $\mathcal C_X^i$ using actual cycles: take
all smooth quasi-projective $T$ and all
$\xi\in\operatorname{CH}^i(T\times X)$. The sum of their
Yoneda maps gives an epimorphism
\[
 \bigoplus_{(T,\xi)}\mathbb Z_{\mathrm{tr}}(T)
                       \longrightarrow\mathcal C_X^i.
\]
Indeed actual Chow sections generate the associated sheaf,
and affine smooth varieties suffice to cover the smooth site.
Thus no assertion that every section of the sheafification
has a global cycle representative is used. The standard
presheaf presentation and its compatibility with
sheafification can also be read in the proof of
Lemma 1.1.6, author pp.~7--8, of the cited paper.

Put $K_T=\ker(\mathbb Z_{\mathrm{tr}}(T)\to
\operatorname{Alb}(T))$, let $\mathcal N$ be the sum of
the images of all $K_T$ in $\mathcal C_X^i$, and let
$\mathcal Q=\mathcal C_X^i/\mathcal N$. The preceding
epimorphism factors as an epimorphism
\[
            \bigoplus_{(T,\xi)}\operatorname{Alb}(T)
                              \longrightarrow\mathcal Q.
\]
The Serre and coproduct properties show that $\mathcal Q$
is one-motivic. Every map from $\mathcal C_X^i$ to a
one-motivic sheaf kills $\mathcal N$, by the same universal
property for each $T$. Consequently
$\mathcal Q=\operatorname{Alb}_{\mathrm{sh}}(\mathcal C_X^i)$,
with its stated unit. This argument uses cokernels and
epimorphisms, and does not require an arbitrary colimit to
be left exact.

Evaluation of etale sheaves at $\operatorname{Spec}\mathbb C$
is exact and commutes with sums: every etale cover of this
point has a section. Sheafification also does not change
the value of the Chow presheaf at this point. Hence
$\mathcal N(\mathbb C)$ is exactly the sum of the images
of the groups $K_T(\mathbb C)$. The latter consists of the
integral zero-cycles on $T$ whose image in its
Serre--Albanese scheme is zero.

We compare this sum with $N^i(X)$, including the issue of
nonproper parameters. Choose a smooth projective
compactification $T\subset\overline T$, separately on
each component. The morphism to the ordinary Albanese of
$\overline T$ factors through $\operatorname{Alb}(T)$.
Thus a zero-cycle $\nu\in K_T(\mathbb C)$ has degree zero
on each component and has zero Albanese class on
$\overline T$. By Chow localization, $\xi$ extends to
$\overline\xi\in\operatorname{CH}^i(\overline T\times X)$.
Its action on this zero-cycle is unchanged:
\[
                       \xi_*(\nu)=\overline\xi_*(\nu).
\]
The right side is one of the defining generators of
$N^i(X)$. Conversely, every defining generator of $N^i(X)$
already uses a smooth projective $T$ and an actual cycle
$\xi$, and is the image of an element of $K_T(\mathbb C)$.
For the smooth projective compactifications,
correspondence actions factor through $\operatorname{CH}_0$;
the argument does not require such a factorization for a
nonproper parameter. This proves
\eqref{eq:E3-sheaf-kernel} integrally.

Finally, the universal property gives
$\pi_0\mathcal E_X^i=\pi_0\mathcal C_X^i$.
Theorem 3.1.4 of the cited paper, author p.~49, identifies
the latter with the sheaf of cycles modulo relative
algebraic equivalence, integrally over $\mathbb C$.
At the point it is
$\operatorname{CH}^i(X)/\operatorname{CH}^i(X)_{\mathrm{alg}}$.
Taking the kernel of the component map, using
$N^i(X)\subset\operatorname{CH}^i(X)_{\mathrm{alg}}$ and
exact evaluation, proves \eqref{eq:E3-sheaf-connected}.
\end{proof}

This theorem alone does not imply finite-dimensionality. The
structure theorem for a general one-motivic sheaf explicitly
allows torsion-free discrete kernels and filtered colimits
of abelian varieties
\citep[Theorem 1.3.10, author p.~23]{cyAyoubBarbieriViale2009}.
For example, if $B/\mathbb C$ is an elliptic curve and
$b\in B(\mathbb C)$ is nontorsion, the etale sheaf
$B/\mathbb Z b$ is connected and one-motivic. Its quotient
map cannot be represented by a morphism to an abelian
variety inducing that point quotient: any such morphism
would kill the Zariski-dense subgroup $\mathbb Zb$, hence
all of $B$, although $B(\mathbb C)/\mathbb Zb\ne0$.
This is a countermodel to a categorical shortcut, not a
claim that it occurs as $P^i(X)$ for a specified variety.

Nor can we apply preservation of compact objects directly
to the Chow internal Hom. The cited paper proves that
$L\operatorname{Alb}$ preserves compact motives with
rational coefficients in Proposition 2.4.7, author p.~42,
but explains on author pp.~46--47 that
\[
 \underline{\operatorname{Hom}}_{\mathrm{eff}}
                   (M(X),\mathbb Z(2)[4])
\]
is nonconstructible for a generic quintic threefold $X$:
its zeroth component quotient is the infinitely generated
codimension-two group of cycles modulo algebraic equivalence.
Neither this nonconstructibility statement nor the sheaf
description above proves that the connected part fails to
be an abelian variety. E.3u below separately refutes the
general comparison with the finite Walker Jacobian by a
specified family detected by full morphic Abel--Jacobi.
There is a further distinction in
\citet[Remark 3.3.6, author p.~52]{cyAyoubBarbieriViale2009}:
after identifying the higher Albanese sheaf with the
Albanese of the Chow sheaf, the authors anticipate its
comparison with Walker's \emph{full morphic} Abel--Jacobi
map; they do not assert that comparison as a theorem.
That map has a potentially much larger target than the
finite-dimensional Walker Jacobian used in E.3m.
The rational comparison with the full morphic target is
also expressly posed as Conjecture 5 in
\citet[p.~2 of the opened author report]{cyAyoub2009Report}.
It is an additional comparison question; we do not use it
as an input to the counterexample or to the identification
\eqref{eq:E3-sheaf-kernel}.

% Full source/proof audits passed by root and frobenius.
\paragraph{E.3t (Proved: the full morphic obstruction to Albanese equivalence).}
Let $X/\mathbb C$ be smooth projective of dimension $n$, let
$0\le i\le n$, and put $r=n-i$. We distinguish the full
morphic Jacobian
\[
 J^{\mathrm{mor}}_r(X)
   =J\bigl(L_rH_{2r+1}(X,\mathbb Z)\bigr)
\]
from the finite-dimensional Walker Jacobian $J_W^i(X)$.
Here $L_rH_m$ is Lawson homology with the inductive mixed
Hodge structure of \citet[Sections 3--5]{cyWalker2007Morphic};
the convention already incorporates the twist making the
relevant pure part have weight $-1$. There are canonical
surjections
\[
 P^i(X)\longrightarrow J^{\mathrm{mor}}_r(X)
       \longrightarrow J_W^i(X)(\mathbb C)
       \longrightarrow J_a^i(X)(\mathbb C).
\]
The last map is an isogeny of abelian varieties. No such
finite-dimensionality assertion applies to the first
morphic target. In particular
\begin{equation}\label{eq:E3-N-killed-morphic}
        N^i(X)\subseteq\ker\Phi^{\mathrm{mor}}_r.
\end{equation}

More precisely, for an abelian variety $A/\mathbb C$ and
$Z\in\operatorname{CH}^i(A\times X)$, let
\[
 \lambda_Z:H_1(A,\mathbb Q)
           \longrightarrow L_rH_{2r+1}(X,\mathbb Q)
\]
be the Lawson cylinder map induced by this correspondence.
Then the composite of $u_Z$ with $P^i(X)\to
J^{\mathrm{mor}}_r(X)$ is zero if and only if
$\lambda_Z=0$. Consequently
\begin{equation}\label{eq:E3-Lawson-necessary}
 u_Z=0\ \Longrightarrow\ \lambda_Z=0
       \ \Longrightarrow\ a({}^tZ)=0.
\end{equation}
The same first implication holds if $u_Z\otimes\mathbb Q=0$.
The next paragraph shows that the second implication cannot
be reversed even for an elliptic parameter.

\begin{proof}
Example 5.5 of \citet[p.~928]{cyWalker2007Morphic} identifies
the morphic Abel--Jacobi map of zero-cycles on every smooth
projective $T$ with its ordinary Albanese map.
Theorem 5.7(2),(4), printed p.~930, says that the morphic
map factors through rational equivalence and is functorial
for all correspondences between smooth projective varieties
with nonnegative dimension shift. A generator of $N^i(X)$
is $W_*\nu$, where
\[
 W\in\operatorname{CH}^i(T\times X),\qquad
 \nu\in\operatorname{CH}_0(T)^0,\qquad
 \operatorname{alb}_T(\nu)=0.
\]
Its dimension shift is $r\ge0$, so the theorem applies and
gives $\Phi^{\mathrm{mor}}_r(W_*\nu)=0$. This proves
\eqref{eq:E3-N-killed-morphic}. Theorem 5.7(1) gives
surjectivity from algebraically trivial cycles to the full
morphic Jacobian. The factorization through the finite
Walker Jacobian is the displayed factorization immediately
before Corollary 5.9, printed p.~932, and that corollary
gives the finite kernel of $J_W^i\to J_a^i$.

We include the point that detects an entire family. Put
$H=L_rH_{2r+1}(X,\mathbb Z)$. It is a countable inductive
limit of finitely generated mixed Hodge structures, so its
underlying group is countable. By Proposition 3.3, printed
pp.~916--917, its Jacobian is
\[
 J(H)=\frac{H_{\mathbb C}}{F^0H_{\mathbb C}+\operatorname{im}H}
\]
in these weights. The correspondence map from $H_1(A)$ is
a morphism of inductive mixed Hodge structures. If its
rationalization is nonzero, its image is a nonzero finite
pure Hodge structure of weight $-1$. Strictness of the
Hodge filtration implies that the induced complex-linear
map
\[
 H_1(A,\mathbb C)/F^0
               \longrightarrow H_{\mathbb C}/F^0
\]
has nonzero image. That image cannot be contained in the
countable image of $H$, or even in its rational span.
Thus the induced Jacobian map has nontorsion image.
Conversely, if the rational map is zero, the integral
map has finite torsion image; the Jacobian of that finite
Hodge structure is zero, so the induced Jacobian map
vanishes. Functoriality and the zero-cycle comparison
identify this Jacobian map with
$a\mapsto\Phi^{\mathrm{mor}}_r(Z_*([a]-[0]))$.
This proves the equivalence and the first implication in
\eqref{eq:E3-Lawson-necessary}, including its rational
version. Composition with the Lawson-to-singular cycle
map gives the ordinary cohomological cylinder. Poincar\'e
duality identifies its transpose with $a({}^tZ)$.
\end{proof}

\paragraph{E.3u (Refuted: the complex converse, with a specified sixfold and elliptic family).}
The implication
\[
                     a({}^tZ)=0\ \Longrightarrow\ u_Z=0
\]
of 14:38--39 is false over $\mathbb C$, even after
rationalizing its conclusion. We give a codimension-four
counterexample on a smooth projective sixfold, with an
elliptic curve as parameter and with an explicitly
specified point whose class modulo $N^4$ is nontorsion.

\emph{The equations and the cycle.}
Let
\[
 Q=\{x_0y_0+x_1y_1+x_2y_2+x_3y_3=0\}
                  \subset\mathbb P^7_{\mathbb Q},
\]
and set
\[
 \Pi_+=\{x_0=x_1=x_2=x_3=0\},\qquad
 \Pi_-=\{y_0=x_1=x_2=x_3=0\},\qquad
 \zeta=[\Pi_+]-[\Pi_-].
\]
Take the degree-$14$ monomials not divisible by $x_0y_0$
as a basis of $H^0(Q,\mathcal O_Q(14))$. Give $x_0^{14}$
coefficient one and every other basis monomial an
independent indeterminate $t_b$. Put
\[
 K=\mathbb Q(t_b),\qquad \Omega=\overline K,
 \qquad F=x_0^{14}+\sum_{b\ne x_0^{14}}t_b b,
 \qquad Y=Q_\Omega\cap\{F=0\}.
\]
This is a specified geometric generic smooth hypersurface,
of dimension five. The intersections with both $\Pi_\pm$
are proper, and we put
\[
 \gamma=[Y\cap\Pi_+]-[Y\cap\Pi_-]
                   \in\operatorname{CH}_2(Y).
\]
Finally set
\[
 E:\ v^2=u^3-u,\qquad o=\text{the point at infinity},
 \qquad X=Y\times E,
\]
and on the ordered product $E_a\times Y\times E_b$ set
\begin{equation}\label{eq:E3-explicit-morphic-family}
 Z=p_Y^*\gamma\cdot p_{ab}^*
                   (\Delta_E-E_a\times\{o\})
             \in\operatorname{CH}^4(E_a\times X).
\end{equation}
For $a\in E(\mathbb C)$ this gives
\[
                 u_Z(a)=\gamma\times([a]-[o])\pmod{N^4(X)}.
\]
All the stated properties hold after every embedding of
the indicated algebraically closed fields into $\mathbb C$.

\emph{The needed strengthened nonvanishing.}
The quadric calculation in C.1i gives
\[
 \deg\zeta^2=-2,
 \qquad h\operatorname{cl}(\zeta)=0,
 \qquad \operatorname{cl}(\zeta)\ne0.
\]
Integral weak Lefschetz identifies $H_4(Y,\mathbb Z)$
with $H_4(Q,\mathbb Z)\simeq\mathbb Z$. The two
surfaces defining $\gamma$ have the same degree, namely
$14$, so
\begin{equation}\label{eq:E3-gamma-lawson}
 \operatorname{cl}(\gamma)=0\text{ in }H_4(Y,\mathbb Z),
 \qquad
 s\gamma\ne0\text{ in }L_1H_4(Y,\mathbb Q).
\end{equation}
Here the second assertion is stronger than Griffiths
nonvanishing, and requires a separate argument.

We use \citet[Theorem 3.4, printed p.~1082, with its proof
on pp.~1095--1097]{cyFriedlander2000}. This is the theorem
behind Walker's Theorem 6.1; the latter's locator ``4.5''
does not refer to the required theorem in the published
Friedlander paper. We spell out why the effective degree
and the geometric generic field are legitimate. For
$S=\mathbb P H^0(Q,\mathcal O_Q(14))$ and every smooth
base $T\to S^\circ$, the Nagel bound verified in C.1i
gives
\begin{equation}\label{eq:E3-effective-Lawson-connectivity}
 H^k(Q\times T,Y_T;\mathbb Q)=0\qquad(k\le10).
\end{equation}
In Nagel's notation we have $c=\mu=5$, $r=0$, and
$m_Q=1$. The two numerical conditions reduce to
$14\ge6$ and $14\ge12$. Thus no unspecified large-degree
threshold is being used.

For completeness, suppose that a nonzero multiple of
$\gamma$ has zero $s$-image. Theorem 3.2 and
Proposition 5.1 of the same Friedlander paper express
such a relation by a projective threefold $W$, allowed
to be singular, a homologically trivial one-cycle
$\delta$ on $W$, and a Chow correspondence from $W$
to the curve space of $Y$. Finite sums are combined
using disjoint unions. These are finite algebraic data.
Spread them on a smooth dominant base $T\to S^\circ$
and shrink so that the proper families and their
resolutions have the required topological strata and
so that the zero homology class of $\delta$ persists.
The finite-type relative Chow construction in
Proposition 5.1 supplies precisely this spread.
We do not assert that coefficients of an arbitrary
complex witness are algebraic over $K$: a smooth
dominant base is sufficient, since
\eqref{eq:E3-effective-Lawson-connectivity} permits
every such base. The smooth locus of that dominant
map is nonempty in characteristic zero.

We recall the decisive cohomology argument, to identify
the only connectivity degrees used. Resolve the
proper family $W$ by $W'$ and resolve
$W'\times_W W'$ by $W''$; shrink $T$ so their
composites to $T$ are smooth. The relative Chow
correspondence has a characteristic class of degree
eight on $W'\times_TY_T$. Connectivity lifts it
uniquely to a class $b$ of degree eight on
$W'\times Q$, and its two pullbacks to
$W''\times Q$ are equal, by the same uniqueness.
The degree-six injectivity from connectivity says
that its action on the lifted cycle $\delta'$ is
the pullback of the nonzero multiple of
$\operatorname{cl}(\zeta)$. On any fibre, proper
homological descent, as used in the proof on p.~1097,
gives
\[
 [\delta']=(p_1)_*\eta-(p_2)_*\eta
                  \quad\text{for some }\eta\in H_2(W'',\mathbb Q),
\]
because its pushforward $[\delta]$ on $W$ is zero.
The projection formula and $p_1^*b=p_2^*b$ make
the action of $b$ on $[\delta']$ zero. This would
give $m\operatorname{cl}(\zeta)=0$ for a nonzero
integer $m$, a contradiction.
The only restriction degrees are eight and six;
relative vanishing through degree nine suffices,
and \eqref{eq:E3-effective-Lawson-connectivity}
gives vanishing through ten. This proves the second
assertion of \eqref{eq:E3-gamma-lawson}.

This argument also explains transfer for the specified
field. A witness over an extension is spread over a
finite-type model of that extension; proper
topological stratification, or its constructible
cohomology formulation, preserves its zero cycle
class after shrinking. The projection to $S$ is
dominant because the coefficients $t_b$ are
algebraically independent. Its smooth locus is
therefore available, and the preceding argument
applies there. It contradicts a witness after any
embedding into $\mathbb C$, without assuming that
an unspecified analytic exceptional locus is
defined over $\overline{\mathbb Q}$.

\emph{The whole cylinder vanishes, while the family survives.}
Equation \eqref{eq:E3-gamma-lawson} implies
$\operatorname{cl}(Z)=0$ integrally by the exterior
product formula, and therefore
\[
 Z_*=0:H_1(E,\mathbb Q)\longrightarrow
                         H_5(X,\mathbb Q).
\]
Equivalently the cohomological cylinder
$H^1(E,\mathbb Q(1))\to H^7(X,\mathbb Q(4))$
is zero, and
$a({}^tZ)=0$ in the manuscript's degree.

On Lawson homology the same family gives
\[
 \lambda_\gamma:H_1(E,\mathbb Q)\longrightarrow
                       L_2H_5(X,\mathbb Q),
              \qquad v\longmapsto\gamma\times v.
\]
This is injective. Indeed choose a class $w$ on the
elliptic curve with nonzero intersection pairing
against a given $v\ne0$. Cap with the pullback of
$w$ and push to $Y$. The product calculation in
\citet[proof of Theorem 6.2, printed pp.~933--934]
{cyWalker2007Morphic} gives
\[
 (p_Y)_*\bigl((\gamma\times v)\cap p_E^*w\bigr)
                      =\langle v,w\rangle\,s\gamma\ne0.
\]
The nondegenerate intersection form on $H_1(E)$
and \eqref{eq:E3-gamma-lawson} prove injectivity.
Strictness of the Hodge filtration now injects the
one-dimensional complex space $H_1(E,\mathbb C)/F^0$
into the Lawson complex quotient used in E.3t.
The integral Lawson group is countable, so the
induced map
\[
 E(\mathbb C)\longrightarrow J^{\mathrm{mor}}_2(X),
           \qquad a\longmapsto
              \Phi^{\mathrm{mor}}_2(\gamma\times([a]-[o]))
\]
has countable kernel and countable inverse image
of the torsion subgroup. In particular it has
nontorsion image. Since it kills $N^4(X)$,
$u_Z\otimes\mathbb Q\ne0$. This proves the
counterexample for the entire family.

We can specify a point witness as well. Adjoin one
further indeterminate $t$ to $\Omega$, choose
$v_t$ with $v_t^2=t^3-t$, put
\[
 \Omega'=\overline{\Omega(t)},\qquad
 e=(t,v_t)\in E(\Omega'),\qquad
 z=\gamma\times([e]-[o])\in\operatorname{CH}^4(X_{\Omega'}).
\]
Then $z$ is nontorsion modulo $N^4$ after every
embedding into $\mathbb C$. Suppose that $mz\in N^4$
for $m\ne0$. An expression by finitely many defining
Albanese-kernel correspondences, including its
rational-equivalence witnesses, spreads over a
smooth dominant finite-type base above an open
of $E_\Omega$. Smooth projective parameter
varieties, their relative Albanese schemes, the
marked zero-cycles and the correspondence cycles
all spread after shrinking. The equation that
each Albanese class is zero is an equation of
sections of its relative abelian scheme and thus
continues to hold. The Chow relation also
continues to hold because its finite rational-
equivalence witnesses have been spread.
Consequently an entire nonempty complex open
subset of $E$ would have morphic image killed by
$m$. This contradicts the countability just
proved. This descent argument permits auxiliary
coefficients transcendental over $\Omega'$; a
dominant parameter base suffices. It proves the
asserted point witness and transfer.

The finite Walker family map is zero, since it
is an algebraic homomorphism from $E$ with zero
map on $H_1$. Thus on this same sixfold
\[
 N^4(X)\subsetneq\ker\psi_W^4
                    =(\mathcal A*\mathcal H)^4(X),
 \qquad
 \ker\bigl(P^4(X)\to J_W^4(X)(\mathbb C)\bigr)
                         \otimes\mathbb Q\ne0.
\]
By E.3o the proposed general equality
$\langle\mathcal H^{*2}\rangle_0=\mathcal A*\mathcal H$
is therefore false. This does not contradict the
proved surface-support cases E.3j, E.3l and E.3q,
nor does it refute the separate possible comparison
between the large one-motivic quotient and the
\emph{full} morphic Jacobian in E.3s. The latter
comparison is an additional question, not needed
to decide the requested converse.\hfill$\square$

The finite certificate
\texttt{code/check\_morphic\_counterexample.py} verifies
the quadric equations and gradients, the two planes,
the self-intersection coefficient and $\zeta^2=-2$,
the $65\,892$ basis monomials, the degree-$14$
connectivity inequalities, the elliptic discriminant
$64$, and the unimodular intersection matrix.
Its output records the exact runtime. The infinite
generic-fibre and Lawson assertions follow from
the proof, not from the bounded computation.


\paragraph{Outcome for E.3 and Lead E4.}
The positive-characteristic direction is proved in E.3e,
with its integral premise retained. E.3u refutes the
complex converse, even rationally, by the displayed
codimension-four family on a sixfold. E.3t proves the stronger
necessary condition that the entire Lawson cylinder
vanish; the example has a zero singular cylinder and
an injective Lawson cylinder. E.3m--p remain exact
equivalences for each fixed target, but their universal
conditions are false: on the specified sixfold
$\langle\mathcal H^{*2}\rangle_0^4(X)
\subsetneq(\mathcal A*\mathcal H)^4(X)$, and the
corresponding quotient has nontorsion classes. The left
side is $N^4(X)$ and needs only surface Albanese kernels
as generators; the right side is exactly $\ker\psi_W$.
E.3j proves that converse for one-cycles whenever zero-cycles are
supported on a curve; E.3l extends it to support on a surface
and identifies $P^{n-1}(X)$ integrally with the Walker Jacobian
group in that case. E.3q adds codimension two under the same
surface-support hypothesis, with the ordinary Jacobian as well.
Applied to the Ottem--Suzuki fourfold,
this combines with E.3h to refute the unqualified integral
biconditional in Lead E4: its family-map assertion is true and
its equality of equivalence relations is false on the same
variety. The rational formulation remains proved by E.3c;
E.3r replaces the existence choices by a displayed generic
matrix and independent Legendre curve, and constructs a
nonzero cycle of order two with the stated failure.
E.3s identifies the manuscript's quotient with the connected
point group of the large one-motivic Chow Albanese sheaf.
Its further comparison with the full morphic Jacobian is
an additional conjectural question, and is not an input
to the decision of E.3 or Lead E4.
The manuscript's degree
$H^{2(n-i)+1}(X)(n-i)$ is the type-correct one; the shorter
$H^{(n-i)+1}$ and the alternate $+3$ in the transcribed reduction
to a curve cannot be used. A curve inside $A$ generates $A$ by
its Albanese; its map on points is not itself a surjection unless
the dimensions permit that. This corrects that literal step of
14:38 while preserving the reduction through its Jacobian.
